\documentclass[11pt]{article}

\usepackage[margin=1in]{geometry}
\usepackage{amsmath,amssymb,amsthm}
\usepackage{enumitem}
\usepackage{xcolor}
\usepackage{hyperref}

\hypersetup{
  colorlinks=true,
  linkcolor=blue!60!black,
  urlcolor=blue!60!black
}

\newcommand{\R}{\mathbb{R}}
\newcommand{\N}{\mathbb{N}}
\newcommand{\F}{\mathcal{F}}
\newcommand{\G}{\mathcal{G}}
\newcommand{\E}{\mathbb{E}}
\newcommand{\Prob}{\mathbb{P}}
\newcommand{\Var}{\operatorname{Var}}
\newcommand{\1}{\mathbf{1}}
\newcommand{\sgn}{\operatorname{sgn}}
\newcommand{\Knowledge}[1]{%
  \par\smallskip\noindent
  \textbf{Knowledge used.}
  {\small\color{blue!60!black}\raggedright #1\par}
  \par\smallskip
}

\title{Chapter 4 Exercises: Section 4.1}
\author{\textsc{Solutions to Rick Durrett's}\\
\textit{Probability: Theory and Examples}}
\date{\today}

\begin{document}
\maketitle

\noindent
These solutions use the local Chapter 4 LLM/Viki knowledge set in
\path{Probability/LLM_Viki_Dataset/Chapter4_Knowledge}. The cited
knowledge IDs refer to entries in
\path{chapter4_knowledge_pieces.jsonl}.

\section*{Exercise 4.1.1}

\noindent\textbf{Question.}
Let $A$ be an event with $\Prob(A)>0$ and let $G\in\G$. Prove the
conditional Bayes formula
\[
  \Prob(G\mid A)
  =
  \frac{\int_G \Prob(A\mid \G)\,d\Prob}
       {\int_\Omega \Prob(A\mid \G)\,d\Prob}.
\]
When $\G$ is generated by a countable partition, identify the usual
Bayes formula as a special case.

\Knowledge{
\path{durrett_ch4_conditional_expectation_definition};
\path{durrett_ch4_partition_conditioning_formula};
\path{durrett_ch4_conditional_expectation_properties}.
}

\noindent\textbf{Solution.}
By definition, $\Prob(A\mid\G)=\E(\1_A\mid\G)$. Since $G\in\G$,
the defining property of conditional expectation gives
\[
  \int_G \Prob(A\mid\G)\,d\Prob
  =
  \int_G \E(\1_A\mid\G)\,d\Prob
  =
  \int_G \1_A\,d\Prob
  =
  \Prob(A\cap G).
\]
Taking $G=\Omega$ gives
\[
  \int_\Omega \Prob(A\mid\G)\,d\Prob=\Prob(A).
\]
Therefore
\[
  \frac{\int_G \Prob(A\mid \G)\,d\Prob}
       {\int_\Omega \Prob(A\mid \G)\,d\Prob}
  =
  \frac{\Prob(A\cap G)}{\Prob(A)}
  =
  \Prob(G\mid A).
\]

If $\G=\sigma(G_1,G_2,\ldots)$ is generated by a countable partition
with $\Prob(G_i)>0$, then $\Prob(A\mid\G)$ is constant on each $G_i$
and equals $\Prob(A\mid G_i)$. Hence, for the atom $G_i$,
\[
  \Prob(G_i\mid A)
  =
  \frac{\Prob(A\mid G_i)\Prob(G_i)}
       {\sum_j \Prob(A\mid G_j)\Prob(G_j)}.
\]

\section*{Exercise 4.1.2}

\noindent\textbf{Question.}
Prove the conditional Chebyshev inequality: for $a>0$,
\[
  \Prob(|X|\ge a\mid\F)\le a^{-2}\E(X^2\mid\F).
\]

\Knowledge{
\path{durrett_ch4_conditional_expectation_definition};
\path{durrett_ch4_conditional_expectation_properties}.
}

\noindent\textbf{Solution.}
Pointwise,
\[
  \1_{\{|X|\ge a\}}\le \frac{X^2}{a^2}.
\]
Taking conditional expectations with respect to $\F$ and using
monotonicity gives
\[
  \Prob(|X|\ge a\mid\F)
  =
  \E(\1_{\{|X|\ge a\}}\mid\F)
  \le
  \E\left(\frac{X^2}{a^2}\mid\F\right)
  =
  a^{-2}\E(X^2\mid\F).
\]

\section*{Exercise 4.1.3}

\noindent\textbf{Question.}
Prove the conditional Cauchy-Schwarz inequality
\[
  \E(XY\mid\G)^2\le \E(X^2\mid\G)\E(Y^2\mid\G).
\]

\Knowledge{
\path{durrett_ch4_conditional_expectation_properties};
\path{durrett_ch4_taking_out_what_is_known};
\path{durrett_ch4_conditional_expectation_l2_projection}.
}

\noindent\textbf{Solution.}
Assume $X,Y\in L^2$. Put
\[
  A=\E(XY\mid\G),\qquad B=\E(Y^2\mid\G),\qquad C=\E(X^2\mid\G).
\]
On the event $\{B=0\}\in\G$, we have
\[
  \int_{\{B=0\}}Y^2\,d\Prob=\int_{\{B=0\}}B\,d\Prob=0,
\]
so $Y=0$ a.s. on $\{B=0\}$. Thus $A=0$ a.s. on $\{B=0\}$.

On $\{B>0\}$, define the $\G$-measurable random variable
$Z=A/B$. Since conditional expectation is positive,
\[
  0\le \E\bigl((X-ZY)^2\mid\G\bigr).
\]
Using linearity and the rule for taking out $\G$-measurable factors,
\[
\begin{aligned}
  0
  &\le
  \E(X^2\mid\G)-2Z\E(XY\mid\G)+Z^2\E(Y^2\mid\G)\\
  &= C-2(A/B)A+(A^2/B^2)B\\
  &= C-\frac{A^2}{B}.
\end{aligned}
\]
Multiplying by $B>0$ gives $A^2\le BC$ on $\{B>0\}$, while the same
inequality is already true on $\{B=0\}$. Hence
\[
  \E(XY\mid\G)^2\le \E(X^2\mid\G)\E(Y^2\mid\G).
\]

\section*{Exercise 4.1.4}

\noindent\textbf{Question.}
Using regular conditional probabilities, prove the conditional Holder
inequality: if $1<p,q<\infty$ and $p^{-1}+q^{-1}=1$, then
\[
  \E(|XY|\mid\G)
  \le
  \E(|X|^p\mid\G)^{1/p}\E(|Y|^q\mid\G)^{1/q}.
\]

\Knowledge{
\path{durrett_ch4_regular_conditional_distribution};
\path{durrett_ch4_conditional_expectation_definition};
\path{durrett_ch4_conditional_expectation_properties}.
}

\noindent\textbf{Solution.}
Let $\mu(\omega,\cdot)$ be a regular conditional distribution of
$(X,Y)$ given $\G$. Then, for almost every $\omega$,
\[
  \E(|XY|\mid\G)(\omega)
  =
  \int_{\R^2}|xy|\,\mu(\omega,d(x,y)).
\]
For those $\omega$, ordinary Holder's inequality applied to the
probability measure $\mu(\omega,\cdot)$ gives
\[
  \int_{\R^2}|xy|\,\mu(\omega,d(x,y))
  \le
  \left(\int_{\R^2}|x|^p\,\mu(\omega,d(x,y))\right)^{1/p}
  \left(\int_{\R^2}|y|^q\,\mu(\omega,d(x,y))\right)^{1/q}.
\]
The regular conditional distribution formula identifies the two factors
as $\E(|X|^p\mid\G)(\omega)^{1/p}$ and
$\E(|Y|^q\mid\G)(\omega)^{1/q}$. Therefore the desired inequality holds
almost surely.

\section*{Exercise 4.1.5}

\noindent\textbf{Question.}
Give an example on $\Omega=\{a,b,c\}$ for which
\[
  \E(\E(X\mid\F_1)\mid\F_2)
  \ne
  \E(\E(X\mid\F_2)\mid\F_1).
\]

\Knowledge{
\path{durrett_ch4_partition_conditioning_formula};
\path{durrett_ch4_tower_property_smaller_sigma_field_wins}.
}

\noindent\textbf{Solution.}
Let $\Prob$ be uniform on $\Omega=\{a,b,c\}$. Define
\[
  \F_1=\sigma(\{a,b\},\{c\}),\qquad
  \F_2=\sigma(\{a,c\},\{b\}),
\]
and let $X=\1_{\{a\}}$.

First compute $\E(X\mid\F_1)$. It is constant on the atoms of $\F_1$:
\[
  \E(X\mid\F_1)=
  \begin{cases}
  1/2, & \omega\in\{a,b\},\\
  0, & \omega=c.
  \end{cases}
\]
Conditioning this on $\F_2$ gives
\[
  \E(\E(X\mid\F_1)\mid\F_2)=
  \begin{cases}
  1/4, & \omega\in\{a,c\},\\
  1/2, & \omega=b.
  \end{cases}
\]

On the other hand,
\[
  \E(X\mid\F_2)=
  \begin{cases}
  1/2, & \omega\in\{a,c\},\\
  0, & \omega=b.
  \end{cases}
\]
Conditioning this on $\F_1$ gives
\[
  \E(\E(X\mid\F_2)\mid\F_1)=
  \begin{cases}
  1/4, & \omega\in\{a,b\},\\
  1/2, & \omega=c.
  \end{cases}
\]
The two random variables differ, for example at $b$ and $c$. This does
not contradict the tower property, because neither $\F_1\subseteq\F_2$
nor $\F_2\subseteq\F_1$.

\section*{Exercise 4.1.6}

\noindent\textbf{Question.}
Suppose $\G\subseteq\F$ and $\E X^2<\infty$. Prove
\[
  \E\{X-\E(X\mid\F)\}^2
  +
  \E\{\E(X\mid\F)-\E(X\mid\G)\}^2
  =
  \E\{X-\E(X\mid\G)\}^2.
\]

\Knowledge{
\path{durrett_ch4_conditional_expectation_l2_projection};
\path{durrett_ch4_taking_out_what_is_known};
\path{durrett_ch4_tower_property_smaller_sigma_field_wins}.
}

\noindent\textbf{Solution.}
Let
\[
  Y=\E(X\mid\F),\qquad Z=\E(X\mid\G).
\]
Since $\G\subseteq\F$, both $Y$ and $Z$ are $\F$-measurable. Also
$Y-Z\in L^2(\F)$. The $L^2$ projection property of conditional
expectation gives
\[
  \E[(X-Y)W]=0\qquad\text{for every }W\in L^2(\F).
\]
Taking $W=Y-Z$, we get
\[
  \E[(X-Y)(Y-Z)]=0.
\]
Now write
\[
  X-Z=(X-Y)+(Y-Z).
\]
Squaring and taking expectations gives
\[
\begin{aligned}
  \E(X-Z)^2
  &=
  \E(X-Y)^2+\E(Y-Z)^2+2\E[(X-Y)(Y-Z)]\\
  &=
  \E(X-Y)^2+\E(Y-Z)^2.
\end{aligned}
\]
Substituting back $Y=\E(X\mid\F)$ and $Z=\E(X\mid\G)$ proves the
identity.

\section*{Exercise 4.1.7}

\noindent\textbf{Question.}
Let
\[
  \Var(X\mid\F)=\E(X^2\mid\F)-\E(X\mid\F)^2.
\]
Prove the variance decomposition
\[
  \Var(X)=\E[\Var(X\mid\F)]+\Var(\E(X\mid\F)).
\]

\Knowledge{
\path{durrett_ch4_conditional_expectation_properties};
\path{durrett_ch4_conditional_expectation_l2_projection}.
}

\noindent\textbf{Solution.}
Let $Y=\E(X\mid\F)$. Taking expectations in the definition of
conditional variance gives
\[
  \E[\Var(X\mid\F)]
  =
  \E[\E(X^2\mid\F)]-\E(Y^2)
  =
  \E X^2-\E Y^2.
\]
Also
\[
  \Var(Y)=\E Y^2-(\E Y)^2.
\]
Since $\E Y=\E[\E(X\mid\F)]=\E X$, adding the two displayed identities
gives
\[
  \E[\Var(X\mid\F)]+\Var(Y)
  =
  \E X^2-(\E X)^2
  =
  \Var(X).
\]

\section*{Exercise 4.1.8}

\noindent\textbf{Question.}
Let $Y_1,Y_2,\ldots$ be i.i.d. with mean $\mu$ and variance $\sigma^2$.
Let $N$ be an independent positive integer-valued random variable with
$\E N^2<\infty$, and define
\[
  X=Y_1+\cdots+Y_N.
\]
Show that
\[
  \Var(X)=\sigma^2\E N+\mu^2\Var(N).
\]

\Knowledge{
\path{durrett_ch4_independent_conditioning_rule};
\path{durrett_ch4_conditional_expectation_properties};
\path{durrett_ch4_tower_property_smaller_sigma_field_wins};
\path{durrett_ch4_conditional_expectation_l2_projection}.
}

\noindent\textbf{Solution.}
Condition on $\sigma(N)$. Given $N=n$, the sum $X$ has mean $n\mu$ and
variance $n\sigma^2$. Thus
\[
  \E(X\mid N)=\mu N
\]
and
\[
  \Var(X\mid N)=\sigma^2 N.
\]
Using the variance decomposition from Exercise 4.1.7,
\[
\begin{aligned}
  \Var(X)
  &=
  \E[\Var(X\mid N)]+\Var(\E(X\mid N))\\
  &=
  \E[\sigma^2N]+\Var(\mu N)\\
  &=
  \sigma^2\E N+\mu^2\Var(N).
\end{aligned}
\]

\section*{Exercise 4.1.9}

\noindent\textbf{Question.}
Suppose $\E(Y\mid\G)=X$ and $\E Y^2=\E X^2<\infty$. Prove that
$X=Y$ almost surely.

\Knowledge{
\path{durrett_ch4_conditional_expectation_l2_projection};
\path{durrett_ch4_taking_out_what_is_known};
\path{durrett_ch4_conditional_expectation_properties}.
}

\noindent\textbf{Solution.}
Since $X=\E(Y\mid\G)$, the random variable $X$ is $\G$-measurable.
Using the rule for taking out what is known,
\[
  \E(XY)
  =
  \E\bigl[X\,\E(Y\mid\G)\bigr]
  =
  \E(X^2).
\]
Therefore
\[
\begin{aligned}
  \E(Y-X)^2
  &=
  \E Y^2-2\E(XY)+\E X^2\\
  &=
  \E Y^2-2\E X^2+\E X^2\\
  &=
  \E Y^2-\E X^2\\
  &=
  0.
\end{aligned}
\]
Hence $Y-X=0$ almost surely.

\section*{Exercise 4.1.10}

\noindent\textbf{Question.}
Assume $\E|Y|<\infty$ and $\E(Y\mid\G)$ has the same distribution as
$Y$. Prove that
\[
  \E(Y\mid\G)=Y\qquad\text{a.s.}
\]

\Knowledge{
\path{durrett_ch4_conditional_expectation_properties};
\path{durrett_ch4_conditional_jensen_contraction};
\path{durrett_ch4_taking_out_what_is_known}.
}

\noindent\textbf{Solution.}
Let $Z=\E(Y\mid\G)$. We first prove a sign lemma. Suppose $U\in L^1$,
$M=\E(U\mid\G)$, and $M$ has the same distribution as $U$. Then
\[
  \E|M|=\E|U|.
\]
Since $\sgn(M)$ is $\G$-measurable and bounded,
\[
  \E|M|
  =
  \E[\sgn(M)M]
  =
  \E[\sgn(M)\E(U\mid\G)]
  =
  \E[\sgn(M)U]
  \le
  \E|U|.
\]
The first and last quantities are equal, so equality must hold in
$\sgn(M)U\le |U|$. Therefore $\sgn(M)U=|U|$ almost surely. In
particular, $U$ and $M$ have the same sign almost surely, with both zero
on the set where $M=0$.

Now apply the lemma to
\[
  U_c=Y-c,\qquad M_c=Z-c=\E(Y-c\mid\G),
\]
where $c\in\mathbb{Q}$. Because $Z$ and $Y$ have the same distribution,
$M_c$ and $U_c$ also have the same distribution. Hence, for each
rational $c$,
\[
  \sgn(Y-c)=\sgn(Z-c)\qquad\text{a.s.}
\]
Outside one null set, this equality holds for every rational $c$.
If $Y(\omega)<Z(\omega)$, choose a rational $c$ strictly between them;
then $Y(\omega)-c<0$ while $Z(\omega)-c>0$, a contradiction. The case
$Z(\omega)<Y(\omega)$ is the same. Hence $Y=Z$ almost surely, i.e.
\[
  \E(Y\mid\G)=Y\qquad\text{a.s.}
\]

\section*{Exercise 4.2.1}

\noindent\textbf{Question.}
Suppose $(X_n)$ is a martingale with respect to a filtration
$(\G_n)$. Let
\[
  \F_n=\sigma(X_1,\ldots,X_n).
\]
Show that $\G_n\supseteq\F_n$ and that $(X_n)$ is also a martingale
with respect to $(\F_n)$.

\Knowledge{
\path{durrett_ch4_martingale_definition};
\path{durrett_ch4_tower_property_smaller_sigma_field_wins};
\path{durrett_ch4_multi_step_martingale_property}.
}

\noindent\textbf{Solution.}
Since $X_i$ is $\G_i$-measurable and $\G_i\subseteq\G_n$ for
$i\le n$, each $X_i$ with $i\le n$ is $\G_n$-measurable. Hence
\[
  \F_n=\sigma(X_1,\ldots,X_n)\subseteq \G_n.
\]
The process is integrable by hypothesis, and $X_n$ is
$\F_n$-measurable by definition. Finally, using the tower property and
$\F_n\subseteq\G_n$,
\[
  \E(X_{n+1}\mid\F_n)
  =
  \E\{\E(X_{n+1}\mid\G_n)\mid\F_n\}
  =
  \E(X_n\mid\F_n)
  =
  X_n.
\]
Thus $(X_n)$ is a martingale with respect to $(\F_n)$.

\section*{Exercise 4.2.2}

\noindent\textbf{Question.}
Give an example of a submartingale $(X_n)$ such that $(X_n^2)$ is a
supermartingale.

\Knowledge{
\path{durrett_ch4_martingale_definition};
\path{durrett_ch4_conditional_expectation_properties}.
}

\noindent\textbf{Solution.}
Take a deterministic process
\[
  X_n=-\frac{1}{n+1},\qquad n\ge 0,
\]
on any probability space with the trivial filtration. Since
$X_{n+1}\ge X_n$, we have
\[
  \E(X_{n+1}\mid\F_n)=X_{n+1}\ge X_n,
\]
so $(X_n)$ is a submartingale. But
\[
  X_{n+1}^2=\frac{1}{(n+2)^2}\le \frac{1}{(n+1)^2}=X_n^2,
\]
so
\[
  \E(X_{n+1}^2\mid\F_n)=X_{n+1}^2\le X_n^2.
\]
Therefore $(X_n^2)$ is a supermartingale.

\section*{Exercise 4.2.3}

\noindent\textbf{Question.}
Show that if $(X_n)$ and $(Y_n)$ are submartingales with respect to
$(\F_n)$, then $(X_n\vee Y_n)$ is also a submartingale.

\Knowledge{
\path{durrett_ch4_martingale_definition};
\path{durrett_ch4_conditional_expectation_properties};
\path{durrett_ch4_convex_transform_submartingale}.
}

\noindent\textbf{Solution.}
The random variable $X_n\vee Y_n$ is $\F_n$-measurable and integrable
because
\[
  |X_n\vee Y_n|\le |X_n|+|Y_n|.
\]
Also,
\[
  X_{n+1}\vee Y_{n+1}\ge X_{n+1},
  \qquad
  X_{n+1}\vee Y_{n+1}\ge Y_{n+1}.
\]
Taking conditional expectations and using monotonicity,
\[
  \E(X_{n+1}\vee Y_{n+1}\mid\F_n)
  \ge \E(X_{n+1}\mid\F_n)\ge X_n,
\]
and similarly
\[
  \E(X_{n+1}\vee Y_{n+1}\mid\F_n)\ge Y_n.
\]
Hence
\[
  \E(X_{n+1}\vee Y_{n+1}\mid\F_n)\ge X_n\vee Y_n,
\]
so $(X_n\vee Y_n)$ is a submartingale.

\section*{Exercise 4.2.4}

\noindent\textbf{Question.}
Let $(X_n)_{n\ge0}$ be a submartingale with
$\sup_n X_n<\infty$ almost surely. Put
\[
  \xi_n=X_n-X_{n-1},\qquad n\ge1,
\]
and suppose
\[
  \E\left(\sup_{n\ge1}\xi_n^+\right)<\infty.
\]
Show that $X_n$ converges almost surely.

\Knowledge{
\path{durrett_ch4_stopped_process_supermartingale};
\path{durrett_ch4_upcrossing_inequality};
\path{durrett_ch4_martingale_convergence_theorem}.
}

\noindent\textbf{Solution.}
Let
\[
  M=\sup_{n\ge1}\xi_n^+.
\]
Then $M\in L^1$ by assumption. Fix a real number $c$ and define the
stopping time
\[
  T_c=\inf\{n\ge0:X_n\ge c\}.
\]
The stopped process $X_{n\wedge T_c}$ is a submartingale. Moreover, on
the event $\{T_c<\infty\}$, if $T_c\ge1$, then
\[
  X_{T_c}=X_{T_c-1}+\xi_{T_c}<c+\xi_{T_c}^+\le c+M,
\]
while if $T_c=0$ the stopped value is fixed at $X_0$. Thus the positive
parts of the stopped process are bounded by an integrable random
variable depending on $c$, $X_0$, and $M$. Hence
\[
  \sup_n \E\bigl(X_{n\wedge T_c}^+\bigr)<\infty.
\]
By the martingale convergence theorem, the stopped process
$X_{n\wedge T_c}$ converges almost surely. Equivalently, for every
rational interval $[a,b]$, the stopped process has only finitely many
upcrossings of $[a,b]$ almost surely.

Since $\sup_n X_n<\infty$ almost surely, the events
\[
  A_c=\{\sup_n X_n<c\},\qquad c\in\mathbb{N},
\]
increase to an event of probability one. On $A_c$, the stopped process
at $T_c$ is just the original process, so the original process has only
finitely many upcrossings of every rational interval $[a,b]$. Therefore
\[
  \liminf_{n\to\infty}X_n=\limsup_{n\to\infty}X_n
\]
almost surely. Thus $X_n$ converges almost surely.

\section*{Exercise 4.2.5}

\noindent\textbf{Question.}
Give an example of a martingale $(X_n)$ such that
\[
  X_n\to -\infty\qquad\text{a.s.}
\]

\Knowledge{
\path{durrett_ch4_martingale_definition};
\path{durrett_ch4_random_walk_martingales}.
}

\noindent\textbf{Solution.}
Let $\xi_1,\xi_2,\ldots$ be independent random variables with
\[
  \Prob(\xi_n=-1)=1-\frac{1}{n^2},
  \qquad
  \Prob(\xi_n=n^2-1)=\frac{1}{n^2}.
\]
Then
\[
  \E\xi_n
  =
  -\left(1-\frac{1}{n^2}\right)
  +(n^2-1)\frac{1}{n^2}
  =
  0.
\]
Define
\[
  X_n=\sum_{k=1}^n \xi_k,\qquad \F_n=\sigma(\xi_1,\ldots,\xi_n).
\]
Then $(X_n)$ is a martingale.

Since
\[
  \sum_{n=1}^\infty \Prob(\xi_n\ne -1)
  =
  \sum_{n=1}^\infty \frac{1}{n^2}<\infty,
\]
the Borel-Cantelli lemma implies that only finitely many of the large
positive jumps occur almost surely. After the last such jump, all
increments are equal to $-1$, so $X_n\to-\infty$ almost surely.

\section*{Exercise 4.2.6}

\noindent\textbf{Question.}
Let $Y_1,Y_2,\ldots$ be nonnegative i.i.d. random variables with
$\E Y_m=1$ and $\Prob(Y_m=1)<1$. Example 4.2.3 shows that
\[
  X_n=\prod_{m=1}^nY_m
\]
is a martingale.
\begin{enumerate}[label=(\roman*)]
\item Use the nonnegative supermartingale convergence theorem and a
contradiction argument to show that $X_n\to0$ almost surely.
\item Use the strong law of large numbers to identify the exponential
rate of decay.
\end{enumerate}

\Knowledge{
\path{durrett_ch4_random_walk_martingales};
\path{durrett_ch4_nonnegative_supermartingale_convergence};
\path{durrett_ch4_martingale_convergence_theorem}.
}

\noindent\textbf{Solution.}
Since $(X_n)$ is a nonnegative martingale, it is also a nonnegative
supermartingale. Therefore $X_n\to X_\infty$ almost surely for some
finite random variable $X_\infty\ge0$.

We claim that $X_\infty=0$ a.s. Since $Y_1$ is not almost surely equal
to $1$, there is an $\varepsilon>0$ such that
\[
  p=\Prob(|Y_1-1|>\varepsilon)>0.
\]
By independence and the second Borel-Cantelli lemma,
$|Y_n-1|>\varepsilon$ occurs infinitely often almost surely. If
$X_\infty>0$, then eventually $X_n>0$ and
\[
  Y_n=\frac{X_n}{X_{n-1}}\to1,
\]
because both $X_n$ and $X_{n-1}$ converge to the same positive limit.
This contradicts the infinitely many deviations of $Y_n$ from $1$.
Thus $X_\infty=0$ a.s.

For the rate, suppose first that $\log Y_1$ is integrable. Then
\[
  \frac{1}{n}\log X_n
  =
  \frac{1}{n}\sum_{m=1}^n \log Y_m
  \to
  c:=\E\log Y_1
  \qquad\text{a.s.}
\]
By Jensen's inequality,
\[
  c=\E\log Y_1<\log \E Y_1=0,
\]
where the inequality is strict because $Y_1$ is not almost surely
constant. If $\Prob(Y_1=0)>0$, then $X_n=0$ eventually almost surely,
and the same conclusion should be read as decay with rate $-\infty$.

\section*{Exercise 4.2.7}

\noindent\textbf{Question.}
Suppose $y_n>-1$ for all $n$ and
\[
  \sum_{n=1}^\infty |y_n|<\infty.
\]
Show that the infinite product
\[
  \prod_{n=1}^\infty (1+y_n)
\]
exists.

\Knowledge{
\path{durrett_ch4_nonnegative_supermartingale_convergence}.
}

\noindent\textbf{Solution.}
Since $\sum_n |y_n|<\infty$, we have $y_n\to0$. Hence, for all large
$n$, $|y_n|\le1/2$. For such $n$,
\[
  |\log(1+y_n)|\le 2|y_n|.
\]
Thus the series
\[
  \sum_{n=1}^\infty \log(1+y_n)
\]
converges absolutely after deleting finitely many initial terms. Since
each factor $1+y_n$ is positive, the partial products satisfy
\[
  \prod_{n=1}^N(1+y_n)
  =
  \exp\left\{\sum_{n=1}^N\log(1+y_n)\right\}.
\]
The exponent converges to a finite limit, so the product converges to a
finite positive limit.

\section*{Exercise 4.2.8}

\noindent\textbf{Question.}
Let $(X_n)$ and $(Y_n)$ be positive integrable processes adapted to
$(\F_n)$. Suppose
\[
  \E(X_{n+1}\mid\F_n)\le (1+Y_n)X_n
\]
and
\[
  \sum_{n=0}^\infty Y_n<\infty\qquad\text{a.s.}
\]
Prove that $X_n$ converges almost surely to a finite limit by finding a
related supermartingale.

\Knowledge{
\path{durrett_ch4_taking_out_what_is_known};
\path{durrett_ch4_nonnegative_supermartingale_convergence};
\path{durrett_ch4_martingale_definition}.
}

\noindent\textbf{Solution.}
Define
\[
  A_0=1,\qquad
  A_n=\prod_{m=0}^{n-1}(1+Y_m),\qquad n\ge1,
\]
and put
\[
  Z_n=\frac{X_n}{A_n}.
\]
Since $A_n$ is $\F_n$-measurable and positive,
\[
\begin{aligned}
  \E(Z_{n+1}\mid\F_n)
  &=
  \E\left(\frac{X_{n+1}}{A_n(1+Y_n)}\Bigm|\F_n\right)\\
  &=
  \frac{1}{A_n(1+Y_n)}\E(X_{n+1}\mid\F_n)\\
  &\le
  \frac{X_n}{A_n}
  =
  Z_n.
\end{aligned}
\]
Thus $(Z_n)$ is a positive supermartingale. By the nonnegative
supermartingale convergence theorem, $Z_n\to Z_\infty$ almost surely for
some finite $Z_\infty$.

By Exercise 4.2.7, the assumption $\sum_nY_n<\infty$ implies
$A_n\to A_\infty$, where $0<A_\infty<\infty$ almost surely. Therefore
\[
  X_n=A_nZ_n\to A_\infty Z_\infty
\]
almost surely, and the limit is finite.

\section*{Exercise 4.2.9}

\noindent\textbf{Question.}
Prove the switching principle. Suppose $(X_n^1)$ and $(X_n^2)$ are
supermartingales with respect to $(\F_n)$, and $N$ is a stopping time
such that
\[
  X_N^1\ge X_N^2.
\]
Show that both
\[
  Y_n=X_n^1\1_{\{N>n\}}+X_n^2\1_{\{N\le n\}}
\]
and
\[
  Z_n=X_n^1\1_{\{N\ge n\}}+X_n^2\1_{\{N<n\}}
\]
are supermartingales.

\Knowledge{
\path{durrett_ch4_martingale_definition};
\path{durrett_ch4_taking_out_what_is_known};
\path{durrett_ch4_stopped_process_supermartingale}.
}

\noindent\textbf{Solution.}
Both processes are adapted and integrable. We first prove the claim for
$Y$. On $\{N\le n\}$,
\[
  Y_n=X_n^2,\qquad
  \E(Y_{n+1}\mid\F_n)\le \E(X_{n+1}^2\mid\F_n)\le X_n^2=Y_n.
\]
On $\{N>n\}$, either $N>n+1$, in which case $Y_{n+1}=X_{n+1}^1$, or
$N=n+1$, in which case
\[
  Y_{n+1}=X_{n+1}^2\le X_{n+1}^1
\]
by the switching assumption at time $N$. Hence on $\{N>n\}$,
$Y_{n+1}\le X_{n+1}^1$, and so
\[
  \E(Y_{n+1}\mid\F_n)\le \E(X_{n+1}^1\mid\F_n)\le X_n^1=Y_n.
\]
Combining the two $\F_n$-measurable cases shows that $Y$ is a
supermartingale.

For $Z$, split into the events $\{N<n\}$, $\{N=n\}$, and $\{N>n\}$.
On $\{N<n\}$, the process follows $X^2$, so the supermartingale
inequality is immediate. On $\{N>n\}$, the process follows $X^1$ at
time $n+1$, so the inequality is also immediate. On $\{N=n\}$,
\[
  Z_n=X_n^1,\qquad Z_{n+1}=X_{n+1}^2,
\]
and therefore
\[
  \E(Z_{n+1}\mid\F_n)
  =
  \E(X_{n+1}^2\mid\F_n)
  \le X_n^2
  \le X_n^1
  =
  Z_n.
\]
Thus $Z$ is a supermartingale.

\section*{Exercise 4.2.10}

\noindent\textbf{Question.}
Prove Dubins' inequality. For every positive supermartingale
$(X_n)_{n\ge0}$, if $U$ is the number of upcrossings of $[a,b]$ with
$0<a<b$, then
\[
  \Prob(U\ge k)
  \le
  \left(\frac{a}{b}\right)^k
  \E\min(X_0/a,1).
\]

\Knowledge{
\path{durrett_ch4_nonnegative_supermartingale_convergence};
\path{durrett_ch4_upcrossing_inequality};
\path{durrett_ch4_martingale_transform_no_gambling_system};
\path{durrett_ch4_stopped_process_supermartingale}.
}

\noindent\textbf{Solution.}
Let
\[
  N_0=-1,
\]
and for $j\ge1$ define the stopping times
\[
  N_{2j-1}=\inf\{m>N_{2j-2}:X_m\le a\},
  \qquad
  N_{2j}=\inf\{m>N_{2j-1}:X_m\ge b\}.
\]
Define $Y_n$ by
\[
  Y_n=1,\qquad 0\le n<N_1,
\]
and, for $j\ge1$,
\[
  Y_n=
  \begin{cases}
  (b/a)^{j-1}X_n/a, & N_{2j-1}\le n<N_{2j},\\
  (b/a)^j, & N_{2j}\le n<N_{2j+1}.
  \end{cases}
\]
At a time $N_{2j-1}$, we switch from the constant value
$(b/a)^{j-1}$ to $(b/a)^{j-1}X_n/a$, and this switch is downward because
$X_{N_{2j-1}}\le a$. At a time $N_{2j}$, we switch from
$(b/a)^{j-1}X_n/a$ to the constant $(b/a)^j$, and this switch is also
downward because $X_{N_{2j}}\ge b$. Since positive constants and
positive multiples of the supermartingale $X$ are supermartingales, the
switching principle implies inductively that
\[
  Z_n^j=Y_{n\wedge N_j}
\]
is a supermartingale for each fixed $j$.

In particular, for $j=2k$,
\[
  \E Y_{n\wedge N_{2k}}\le \E Y_0.
\]
From the definition,
\[
  Y_0=\min(X_0/a,1).
\]
Letting $n\to\infty$ and using Fatou's lemma,
\[
  \E\left[\liminf_{n\to\infty}Y_{n\wedge N_{2k}}\right]
  \le
  \E\min(X_0/a,1).
\]
On the event $\{U\ge k\}$, the stopping time $N_{2k}$ is finite, and
\[
  \liminf_{n\to\infty}Y_{n\wedge N_{2k}}
  =
  Y_{N_{2k}}
  \ge
  (b/a)^k.
\]
Therefore
\[
  (b/a)^k\Prob(U\ge k)
  \le
  \E\min(X_0/a,1),
\]
which is equivalent to
\[
  \Prob(U\ge k)
  \le
  \left(\frac{a}{b}\right)^k
  \E\min(X_0/a,1).
\]

\section*{Exercise 4.3.1}

\noindent\textbf{Question.}
Give a bounded process with bounded increments that visits each of
$-1,0,1$ infinitely often almost surely. Explain why this shows that
bounded increments alone are not enough in Theorem 4.3.1.

\Knowledge{
\path{durrett_ch4_bounded_increment_convergence};
\path{durrett_ch4_martingale_convergence_theorem}.
}

\noindent\textbf{Solution.}
Take the deterministic sequence
\[
  X_{3k}=-1,\qquad X_{3k+1}=0,\qquad X_{3k+2}=1,\qquad k\ge0.
\]
Then $\sup_n |X_n|\le1$ and $\sup_n |X_{n+1}-X_n|\le2$. Each of
$-1,0,1$ occurs infinitely often, so
\[
  \Prob(X_n=a\ \text{i.o.})=1,\qquad a\in\{-1,0,1\}.
\]
The sequence does not converge and does not oscillate between
$+\infty$ and $-\infty$. Thus the martingale hypothesis in Theorem
4.3.1 is essential; bounded increments by themselves do not force the
theorem's dichotomy. A literal bounded martingale example with these
properties cannot exist, since bounded martingales converge almost
surely.

\section*{Exercise 4.3.2}

\noindent\textbf{Question.}
Construct a finite-state Markov chain on $\{-1,0,1\}$ such that
\[
  \Prob(X_n=0)\to 1-2p,\qquad
  \Prob(X_n=-1)\to p,\qquad
  \Prob(X_n=1)\to p
\]
for some chosen $p\in(0,1/2)$, while the chain does not converge almost
surely.

\Knowledge{
\path{durrett_ch4_bounded_increment_convergence};
\path{durrett_ch4_martingale_convergence_theorem}.
}

\noindent\textbf{Solution.}
Choose, for instance, $p=1/4$. Let $(X_n)$ be independent with
\[
  \Prob(X_n=-1)=1/4,\qquad
  \Prob(X_n=0)=1/2,\qquad
  \Prob(X_n=1)=1/4
\]
for every $n$. This is a Markov chain whose transition probabilities do
not depend on the present state: every row of the transition matrix is
$(1/4,1/2,1/4)$.

The distribution is constant in $n$, so it converges to the desired
limits with $p=1/4$. Since the events $\{X_n=a\}$ are independent and
have positive probability for each $a\in\{-1,0,1\}$, the second
Borel-Cantelli lemma implies that each state is visited infinitely often
almost surely. Hence $X_n$ does not converge almost surely, and in fact
does not converge in probability. This illustrates that convergence in
distribution alone does not imply almost sure convergence.

\section*{Exercise 4.3.3}

\noindent\textbf{Question.}
Let $(X_n)$ and $(Y_n)$ be positive integrable adapted processes.
Suppose
\[
  \E(X_{n+1}\mid\F_n)\le X_n+Y_n
\]
and
\[
  \sum_{n=1}^\infty Y_n<\infty\qquad\text{a.s.}
\]
Prove that $X_n$ converges almost surely to a finite limit.

\Knowledge{
\path{durrett_ch4_martingale_definition};
\path{durrett_ch4_stopped_process_supermartingale};
\path{durrett_ch4_nonnegative_supermartingale_convergence}.
}

\noindent\textbf{Solution.}
Let
\[
  S_n=\sum_{m=1}^nY_m,\qquad Z_n=X_n-S_n.
\]
This indexing puts $Y_{n+1}$ outside $\F_n$, so use the predictable
partial sum instead:
\[
  T_n=\sum_{m=0}^{n-1}Y_m,\qquad Z_n=X_n-T_n,
\]
where the hypothesis is read as
$\E(X_{n+1}\mid\F_n)\le X_n+Y_n$. Then
\[
  \E(Z_{n+1}\mid\F_n)
  =
  \E(X_{n+1}\mid\F_n)-T_n-Y_n
  \le
  X_n-T_n
  =
  Z_n.
\]
Thus $Z_n$ is a supermartingale.

Fix $M<\infty$ and stop just before the partial sum of the $Y$'s exceeds
$M$:
\[
  N_M=\inf\{n:T_{n+1}>M\}.
\]
For $n\le N_M$, $T_n\le M$, and since $X_n\ge0$,
\[
  Z_{n\wedge N_M}=X_{n\wedge N_M}-T_{n\wedge N_M}\ge -M.
\]
Hence $Z_{n\wedge N_M}+M$ is a nonnegative supermartingale, so it
converges almost surely. On the event $\{T_\infty\le M\}$, this stopped
process is eventually the original $Z_n$, so $Z_n$ converges there.
Since $T_\infty<\infty$ almost surely, the events $\{T_\infty\le M\}$
increase to probability one. Therefore $Z_n$ converges almost surely.
Also $T_n\to T_\infty<\infty$, so
\[
  X_n=Z_n+T_n
\]
converges almost surely to a finite limit.

\section*{Exercise 4.3.4}

\noindent\textbf{Question.}
Let $p_m\in[0,1)$. Show that
\[
  \prod_{m=1}^\infty(1-p_m)=0
  \qquad\text{if and only if}\qquad
  \sum_{m=1}^\infty p_m=\infty.
\]

\Knowledge{
\path{durrett_ch4_conditional_borel_cantelli};
\path{durrett_ch4_nonnegative_supermartingale_convergence}.
}

\noindent\textbf{Solution.}
If $\sum_m p_m<\infty$, then $p_m\to0$. For all large $m$,
$p_m\le1/2$, and
\[
  \log(1-p_m)\ge -2p_m.
\]
Thus $\sum_m\log(1-p_m)$ converges to a finite real number after
discarding finitely many initial terms. Since each initial factor is
strictly positive, the product is strictly positive.

Conversely, if $\sum_m p_m=\infty$, let $A_m$ be independent events with
$\Prob(A_m)=p_m$. By the second Borel-Cantelli lemma,
$\Prob(A_m\ \text{i.o.})=1$. In particular, the probability that none of
the events ever occurs is zero. But
\[
  \Prob(A_m^c\ \text{for all }m)
  =
  \prod_{m=1}^\infty(1-p_m).
\]
Hence the product is zero.

\section*{Exercise 4.3.5}

\noindent\textbf{Question.}
Show that if
\[
  \sum_{n=2}^\infty
  \Prob\left(A_n\mid\bigcap_{m=1}^{n-1}A_m^c\right)
  =
  \infty,
\]
then
\[
  \Prob\left(\bigcap_{m=1}^\infty A_m^c\right)=0.
\]

\Knowledge{
\path{durrett_ch4_conditional_borel_cantelli};
\path{durrett_ch4_conditional_expectation_definition}.
}

\noindent\textbf{Solution.}
Let
\[
  B_n=\bigcap_{m=1}^{n-1}A_m^c,\qquad
  q_n=\Prob(A_n\mid B_n),
\]
whenever $\Prob(B_n)>0$. If $\Prob(B_n)=0$ for some $n$, then the
conclusion is immediate. Otherwise,
\[
  \Prob(B_{n+1})=\Prob(B_n)\Prob(A_n^c\mid B_n)
  =\Prob(B_n)(1-q_n).
\]
Iterating gives
\[
  \Prob\left(\bigcap_{m=1}^N A_m^c\right)
  =
  \Prob(A_1^c)\prod_{n=2}^N(1-q_n).
\]
By Exercise 4.3.4, the infinite product is zero because
$\sum_nq_n=\infty$. Letting $N\to\infty$ gives
\[
  \Prob\left(\bigcap_{m=1}^\infty A_m^c\right)=0.
\]

\section*{Exercise 4.3.6}

\noindent\textbf{Question.}
For the partition approximation in Example 4.3.7, verify directly that
$X_n$ is a martingale. If the condition $\mu_n\ll\nu_n$ is dropped and
$X_n$ is set to $0$ on atoms with $\nu$-mass zero, show that $X_n$ is a
supermartingale.

\Knowledge{
\path{durrett_ch4_radon_nikodym_derivative_martingale};
\path{durrett_ch4_conditional_expectation_definition};
\path{durrett_ch4_partition_conditioning_formula}.
}

\noindent\textbf{Solution.}
Let $\F_n$ be generated by a finite or countable partition
$\{I_{k,n}\}_k$, with the $(n+1)$st partition refining the $n$th. In the
absolutely continuous case,
\[
  X_n=\frac{\mu(I_{k,n})}{\nu(I_{k,n})}
  \qquad\text{on } I_{k,n}.
\]
Fix an atom $I=I_{k,n}$ with $\nu(I)>0$. Since the next partition
refines the present one,
\[
  \int_I X_{n+1}\,d\nu
  =
  \sum_{J\subseteq I}\int_J X_{n+1}\,d\nu
  =
  \sum_{J\subseteq I}\mu(J)
  =
  \mu(I)
  =
  \int_I X_n\,d\nu.
\]
This verifies the conditional expectation identity on every atom of
$\F_n$, so $\E_\nu(X_{n+1}\mid\F_n)=X_n$.

If $\mu_n\ll\nu_n$ is not assumed and we define $X_n=0$ on atoms with
$\nu$-mass zero, then for an atom $I$ with $\nu(I)>0$,
\[
  \int_I X_{n+1}\,d\nu
  =
  \sum_{\substack{J\subseteq I\\ \nu(J)>0}}\mu(J)
  \le
  \mu(I)
  =
  \int_I X_n\,d\nu.
\]
On atoms with $\nu(I)=0$, both integrals with respect to $\nu$ are zero.
Thus
\[
  \E_\nu(X_{n+1}\mid\F_n)\le X_n,
\]
so $(X_n)$ is a supermartingale.

\section*{Exercise 4.3.7}

\noindent\textbf{Question.}
Use Theorem 4.3.5 and the partition construction of Example 4.3.7 to
prove the Radon-Nikodym theorem when $\F$ is countably generated.

\Knowledge{
\path{durrett_ch4_radon_nikodym_derivative_martingale};
\path{durrett_ch4_nonnegative_supermartingale_convergence};
\path{durrett_ch4_conditional_expectation_definition}.
}

\noindent\textbf{Solution.}
First suppose $\mu$ and $\nu$ are finite and $\mu\ll\nu$. Let
$\F=\sigma(A_1,A_2,\ldots)$ and set
\[
  \F_n=\sigma(A_1,\ldots,A_n).
\]
Each $\F_n$ has finitely many atoms. On each atom $I$ of $\F_n$, define
\[
  X_n=\frac{\mu(I)}{\nu(I)},
\]
with the value irrelevant when $\nu(I)=0$ because $\mu\ll\nu$. By
Exercise 4.3.6, $(X_n)$ is the Radon-Nikodym derivative martingale for
the finite sigma-fields. Theorem 4.3.5 gives a limit $g$ such that
\[
  \mu(A)=\int_A g\,d\nu+\mu(A\cap\{g=\infty\}).
\]
But $\nu(g=\infty)=0$, and $\mu\ll\nu$, so
$\mu(\{g=\infty\})=0$. Hence
\[
  \mu(A)=\int_A g\,d\nu,\qquad A\in\F.
\]
This is the Radon-Nikodym derivative.

For sigma-finite $\mu$ and $\nu$, decompose $\Omega$ into disjoint
measurable sets $E_j$ on which both measures are finite. Apply the
finite result to the restrictions of $\mu$ and $\nu$ to each $E_j$,
obtaining densities $g_j$ on $E_j$. Then
\[
  g=\sum_j g_j\1_{E_j}
\]
satisfies $\mu(A)=\int_A g\,d\nu$ for all $A\in\F$.

\section*{Exercise 4.3.8}

\noindent\textbf{Question.}
For Bernoulli product measures with coordinate probabilities
\[
  \mu(\xi_n=1)=\alpha_n,\qquad \nu(\xi_n=1)=\beta_n,
\]
find the absolute-continuity criterion from Theorem 4.3.8. Then show
that if $0<\varepsilon\le\alpha_n,\beta_n\le1-\varepsilon<1$, the
criterion is equivalent to
\[
  \sum_n(\alpha_n-\beta_n)^2<\infty.
\]

\Knowledge{
\path{durrett_ch4_radon_nikodym_derivative_martingale};
\path{durrett_ch4_conditional_jensen_contraction}.
}

\noindent\textbf{Solution.}
For one coordinate, the Hellinger integral in Theorem 4.3.8 is
\[
  h_n
  =
  \sqrt{\alpha_n\beta_n}
  +
  \sqrt{(1-\alpha_n)(1-\beta_n)}.
\]
Thus, assuming the one-coordinate absolute-continuity condition
($\beta_n=0$ implies $\alpha_n=0$, and $\beta_n=1$ implies
$\alpha_n=1$), Theorem 4.3.8 says
\[
  \mu\ll\nu
  \qquad\Longleftrightarrow\qquad
  \prod_{n=1}^\infty h_n>0.
\]

Now assume the uniform bounds away from $0$ and $1$. Since
\[
  1-h_n
  =
  \frac12\left[
  \bigl(\sqrt{\alpha_n}-\sqrt{\beta_n}\bigr)^2
  +
  \bigl(\sqrt{1-\alpha_n}-\sqrt{1-\beta_n}\bigr)^2
  \right],
\]
and the functions $x\mapsto\sqrt{x}$ and
$x\mapsto\sqrt{1-x}$ have derivatives bounded above and below away from
$0$ on $[\varepsilon,1-\varepsilon]$, there are constants
$0<c<C<\infty$ such that
\[
  c(\alpha_n-\beta_n)^2
  \le
  1-h_n
  \le
  C(\alpha_n-\beta_n)^2.
\]
Since $0<h_n\le1$, the product $\prod_nh_n$ is positive if and only if
$\sum_n(1-h_n)<\infty$. Therefore, under the uniform bounds,
\[
  \mu\ll\nu
  \qquad\Longleftrightarrow\qquad
  \sum_n(\alpha_n-\beta_n)^2<\infty.
\]

\section*{Exercise 4.3.9}

\noindent\textbf{Question.}
In the Bernoulli product setting, show that if
\[
  \sum_n\alpha_n<\infty
  \qquad\text{and}\qquad
  \sum_n\beta_n=\infty,
\]
then $\mu\perp\nu$. Explain why this shows that the square-summability
condition in Exercise 4.3.8 is not sufficient without uniform bounds.

\Knowledge{
\path{durrett_ch4_conditional_borel_cantelli};
\path{durrett_ch4_radon_nikodym_derivative_martingale}.
}

\noindent\textbf{Solution.}
Let
\[
  A=\{\xi_n=1\ \text{infinitely often}\}.
\]
Under $\mu$, the events $\{\xi_n=1\}$ are independent and have
probabilities $\alpha_n$. Since $\sum_n\alpha_n<\infty$, the first
Borel-Cantelli lemma gives $\mu(A)=0$.

Under $\nu$, the same events have probabilities $\beta_n$, and
$\sum_n\beta_n=\infty$. By the second Borel-Cantelli lemma,
$\nu(A)=1$. Hence $\mu$ and $\nu$ are concentrated on disjoint
measurable sets, so $\mu\perp\nu$.

For example, $\alpha_n=n^{-2}$ and $\beta_n=n^{-1}$ satisfy
\[
  \sum_n(\alpha_n-\beta_n)^2<\infty,
\]
but the preceding argument shows that the product measures are singular.

\section*{Exercise 4.3.10}

\noindent\textbf{Question.}
Assume $0<\alpha_n,\beta_n<1$ for all $n$. Show that
\[
  \sum_n|\alpha_n-\beta_n|<\infty
\]
is sufficient for $\mu\ll\nu$.

\Knowledge{
\path{durrett_ch4_radon_nikodym_derivative_martingale};
\path{durrett_ch4_conditional_jensen_contraction}.
}

\noindent\textbf{Solution.}
Since $0<\alpha_n,\beta_n<1$, each one-coordinate Bernoulli law for
$\mu$ is absolutely continuous with respect to the corresponding one for
$\nu$. Let
\[
  h_n=
  \sqrt{\alpha_n\beta_n}
  +
  \sqrt{(1-\alpha_n)(1-\beta_n)}.
\]
For Bernoulli measures,
\[
  1-h_n
  =
  \frac12\left[
  \bigl(\sqrt{\alpha_n}-\sqrt{\beta_n}\bigr)^2
  +
  \bigl(\sqrt{1-\alpha_n}-\sqrt{1-\beta_n}\bigr)^2
  \right].
\]
The Hellinger distance is bounded by total variation, and here total
variation is $|\alpha_n-\beta_n|$. Thus
\[
  1-h_n\le |\alpha_n-\beta_n|.
\]
Therefore $\sum_n(1-h_n)<\infty$. Since $0<h_n\le1$, this implies
\[
  \prod_n h_n>0.
\]
Theorem 4.3.8 now gives $\mu\ll\nu$.

\section*{Exercise 4.3.11}

\noindent\textbf{Question.}
Let $Z_n$ be a supercritical branching process with mean $\mu>1$ and
let
\[
  W=\lim_{n\to\infty}Z_n/\mu^n.
\]
Show that if $\Prob(W=0)<1$, then $\Prob(W=0)=\rho$, the extinction
probability. Conclude that
\[
  \{W>0\}=\{Z_n>0\ \text{for all }n\}\qquad\text{a.s.}
\]

\Knowledge{
\path{durrett_ch4_branching_process_normalized_martingale};
\path{durrett_ch4_nonnegative_supermartingale_convergence}.
}

\noindent\textbf{Solution.}
Let
\[
  q=\Prob(W=0\mid Z_0=1).
\]
Condition on the first generation. If $Z_1=k$, then the future process
is the sum of $k$ independent descendant branching processes. The limit
$W$ is zero exactly when all $k$ descendant limits are zero. Hence
\[
  q=\sum_{k=0}^\infty p_k q^k=\phi(q).
\]
Thus $q$ is a fixed point of the generating function. In the
supercritical case, the fixed points in $[0,1]$ are the extinction
probability $\rho$ and $1$. Since $\Prob(W=0)<1$, we have $q<1$, and
therefore $q=\rho$.

Extinction implies $W=0$, so
\[
  \{W>0\}\subseteq\{Z_n>0\ \text{for all }n\}.
\]
Both events have probability $1-\rho$, so they are equal almost surely.

\section*{Exercise 4.3.12}

\noindent\textbf{Question.}
Let $(Z_n)$ be a branching process with offspring generating function
\[
  \phi(\theta)=\sum_{k=0}^\infty p_k\theta^k.
\]
Suppose $\rho<1$ and $\phi(\rho)=\rho$. Show that $\rho^{Z_n}$ is a
martingale and conclude that, starting from $Z_0=x$,
\[
  \Prob(Z_n=0\ \text{for some }n\ge1)=\rho^x.
\]

\Knowledge{
\path{durrett_ch4_branching_process_normalized_martingale};
\path{durrett_ch4_nonnegative_supermartingale_convergence};
\path{durrett_ch4_tower_property_smaller_sigma_field_wins}.
}

\noindent\textbf{Solution.}
Given $\F_n$ and $Z_n=k$, the next generation is the sum of $k$
independent offspring variables. Therefore
\[
  \E(\rho^{Z_{n+1}}\mid\F_n)
  =
  \phi(\rho)^{Z_n}
  =
  \rho^{Z_n}.
\]
Thus $(\rho^{Z_n})$ is a bounded martingale.

Let
\[
  T=\inf\{n\ge1:Z_n=0\}.
\]
Since the martingale is bounded, $\E_x\rho^{Z_n}=\rho^x$ for every
$n$. Also $\rho^{Z_n}\to\1_{\{T<\infty\}}$: on extinction it is
eventually $1$, while on survival the supercritical branching process
has $Z_n\to\infty$ and hence $\rho^{Z_n}\to0$. Bounded convergence gives
\[
  \rho^x
  =
  \lim_{n\to\infty}\E_x\rho^{Z_n}
  =
  \Prob_x(T<\infty).
\]

\section*{Exercise 4.3.13}

\noindent\textbf{Question.}
In the Galton-Watson family-name model, each family has exactly three
children and each child is male with probability $1/2$. Thus
\[
  p_0=1/8,\quad p_1=3/8,\quad p_2=3/8,\quad p_3=1/8.
\]
Compute the extinction probability when $Z_0=1$.

\Knowledge{
\path{durrett_ch4_branching_process_normalized_martingale}.
}

\noindent\textbf{Solution.}
The offspring generating function is
\[
  \phi(s)
  =
  \frac18+\frac38s+\frac38s^2+\frac18s^3
  =
  \left(\frac{1+s}{2}\right)^3.
\]
The extinction probability $\rho$ is the smallest solution in $[0,1]$
of
\[
  \rho=\phi(\rho)=\left(\frac{1+\rho}{2}\right)^3.
\]
Equivalently,
\[
  8\rho=(1+\rho)^3,
\]
so
\[
  \rho^3+3\rho^2-5\rho+1=0.
\]
Factoring,
\[
  \rho^3+3\rho^2-5\rho+1
  =
  (\rho-1)(\rho^2+4\rho-1).
\]
The root in $[0,1)$ is
\[
  \rho=-2+\sqrt{5}.
\]
Therefore the probability that the family name dies out, starting from
one male, is
\[
  \boxed{\sqrt{5}-2}.
\]

\section*{Exercise 4.4.1}

\noindent\textbf{Question.}
Let $(X_n)$ be a submartingale and let $N$ be a stopping time with
$\Prob(N\le k)=1$. Show that if $j\le k$, then
\[
  \E(X_j;N=j)\le \E(X_k;N=j),
\]
and sum over $j$ to obtain another proof of $\E X_N\le \E X_k$.

\Knowledge{
\path{durrett_ch4_bounded_stopping_submartingale_expectation};
\path{durrett_ch4_multi_step_martingale_property};
\path{durrett_ch4_tower_property_smaller_sigma_field_wins}.
}

\noindent\textbf{Solution.}
Since $N$ is a stopping time, $\{N=j\}\in\F_j$. For $j\le k$,
the submartingale property over several steps gives
\[
  \E(X_k\mid\F_j)\ge X_j.
\]
Multiplying by $\1_{\{N=j\}}$ and taking expectations,
\[
  \E(X_k;N=j)
  =
  \E\{\1_{\{N=j\}}\E(X_k\mid\F_j)\}
  \ge
  \E(X_j;N=j).
\]
Since $N\le k$ almost surely,
\[
  \E X_N
  =
  \sum_{j=0}^k \E(X_j;N=j)
  \le
  \sum_{j=0}^k \E(X_k;N=j)
  =
  \E X_k.
\]

\section*{Exercise 4.4.2}

\noindent\textbf{Question.}
Generalize Theorem 4.4.1 as follows: if $(X_n)$ is a submartingale and
$M\le N$ are stopping times with $\Prob(N\le k)=1$, prove that
\[
  \E X_M\le \E X_N.
\]

\Knowledge{
\path{durrett_ch4_bounded_stopping_submartingale_expectation};
\path{durrett_ch4_stopped_process_supermartingale};
\path{durrett_ch4_martingale_transform_no_gambling_system}.
}

\noindent\textbf{Solution.}
Define the predictable process
\[
  H_m=\1_{\{M<m\le N\}},\qquad 1\le m\le k.
\]
Indeed,
\[
  \{M<m\le N\}=\{M\le m-1\}\cap\{N\ge m\}\in\F_{m-1}.
\]
Since $H_m\ge0$ and $X$ is a submartingale, the martingale-transform
principle says that
\[
  (H\cdot X)_n=\sum_{m=1}^nH_m(X_m-X_{m-1})
\]
is a submartingale starting from $0$. At time $k$,
\[
  (H\cdot X)_k=X_N-X_M,
\]
because the sum includes exactly the increments after $M$ and up to
$N$. Therefore
\[
  \E X_N-\E X_M=\E(H\cdot X)_k\ge0,
\]
which proves $\E X_M\le \E X_N$.

\section*{Exercise 4.4.3}

\noindent\textbf{Question.}
Suppose $M\le N$ are stopping times. If $A\in\F_M$, show that
\[
  L=
  \begin{cases}
  M, & \text{on } A,\\
  N, & \text{on } A^c
  \end{cases}
\]
is a stopping time.

\Knowledge{
\path{durrett_ch4_stopped_process_supermartingale};
\path{durrett_ch4_bounded_stopping_submartingale_expectation}.
}

\noindent\textbf{Solution.}
For each $n$,
\[
  \{L\le n\}
  =
  \bigl(A\cap\{M\le n\}\bigr)
  \cup
  \bigl(A^c\cap\{N\le n\}\bigr).
\]
Since $A\in\F_M$, the event $A\cap\{M\le n\}$ belongs to $\F_n$.
Also, because $M\le N$, on $\{N\le n\}$ we have $M\le n$, and the
definition of $A\in\F_M$ implies
\[
  A^c\cap\{N\le n\}
  =
  A^c\cap\{M\le n\}\cap\{N\le n\}\in\F_n.
\]
Therefore $\{L\le n\}\in\F_n$ for every $n$, so $L$ is a stopping time.

\section*{Exercise 4.4.4}

\noindent\textbf{Question.}
Use Exercise 4.4.3 to strengthen Exercise 4.4.2 to
\[
  X_M\le \E(X_N\mid\F_M)
\]
for bounded stopping times $M\le N$.

\Knowledge{
\path{durrett_ch4_bounded_stopping_submartingale_expectation};
\path{durrett_ch4_conditional_expectation_definition};
\path{durrett_ch4_stopped_process_supermartingale}.
}

\noindent\textbf{Solution.}
Let $A\in\F_M$ and define $L=M$ on $A$ and $L=N$ on $A^c$. By Exercise
4.4.3, $L$ is a stopping time and $M\le L\le N$. Exercise 4.4.2 gives
\[
  \E X_M\le \E X_L\le \E X_N.
\]
Since
\[
  X_L=X_M\1_A+X_N\1_{A^c},
\]
the inequality $\E X_L\le \E X_N$ gives
\[
  \E(X_M;A)+\E(X_N;A^c)\le \E(X_N;A)+\E(X_N;A^c).
\]
Thus
\[
  \E(X_M;A)\le \E(X_N;A)
  \qquad\text{for every }A\in\F_M.
\]
By the defining property of conditional expectation, this is equivalent
to
\[
  X_M\le \E(X_N\mid\F_M)\qquad\text{a.s.}
\]

\section*{Exercise 4.4.5}

\noindent\textbf{Question.}
Prove the following conditional-variance identity. If
$\F\subseteq\G$, then
\[
  \E\left(\E(Y\mid\G)-\E(Y\mid\F)\right)^2
  =
  \E\left(\E(Y\mid\G)\right)^2
  -
  \E\left(\E(Y\mid\F)\right)^2.
\]

\Knowledge{
\path{durrett_ch4_conditional_expectation_l2_projection};
\path{durrett_ch4_tower_property_smaller_sigma_field_wins};
\path{durrett_ch4_taking_out_what_is_known}.
}

\noindent\textbf{Solution.}
Let
\[
  U=\E(Y\mid\G),\qquad V=\E(Y\mid\F).
\]
Since $\F\subseteq\G$, the tower property gives
\[
  \E(U\mid\F)=\E(\E(Y\mid\G)\mid\F)=\E(Y\mid\F)=V.
\]
Therefore
\[
  \E(UV)
  =
  \E\{V\E(U\mid\F)\}
  =
  \E(V^2).
\]
Expanding the square,
\[
\begin{aligned}
  \E(U-V)^2
  &=
  \E U^2-2\E(UV)+\E V^2\\
  &=
  \E U^2-\E V^2,
\end{aligned}
\]
which is the desired identity.

\section*{Exercise 4.4.6}

\noindent\textbf{Question.}
Let $S_n=\xi_1+\cdots+\xi_n$, where the $\xi_m$ are independent with
$\E\xi_m=0$, $\sigma_m^2=\E\xi_m^2$, and $|\xi_m|\le K$. Put
\[
  s_n^2=\sum_{m=1}^n\sigma_m^2=\Var(S_n).
\]
Use the martingale $S_n^2-s_n^2$ and Theorem 4.4.1 to prove
\[
  \Prob\left(\max_{1\le m\le n}|S_m|\le x\right)
  \le
  \frac{(x+K)^2}{\Var(S_n)}.
\]

\Knowledge{
\path{durrett_ch4_random_walk_martingales};
\path{durrett_ch4_bounded_stopping_submartingale_expectation};
\path{durrett_ch4_conditional_variance_formula_martingales}.
}

\noindent\textbf{Solution.}
Let
\[
  T=\inf\{m\ge1:|S_m|>x\}.
\]
The process $S_m^2-s_m^2$ is a martingale. Applying bounded optional
sampling to $T\wedge n$ gives
\[
  \E S_{T\wedge n}^2=\E s_{T\wedge n}^2.
\]
On the event
\[
  A=\left\{\max_{1\le m\le n}|S_m|\le x\right\},
\]
we have $T>n$, so $s_{T\wedge n}^2=s_n^2$. Hence
\[
  \E s_{T\wedge n}^2\ge s_n^2\Prob(A).
\]
On the other hand, if $T\le n$, then the bounded increment assumption
gives $|S_T|\le x+K$, while if $T>n$, $|S_n|\le x$. Thus
\[
  |S_{T\wedge n}|\le x+K
  \qquad\text{always}.
\]
Therefore
\[
  s_n^2\Prob(A)\le \E S_{T\wedge n}^2\le (x+K)^2.
\]
Since $s_n^2=\Var(S_n)$, the result follows.

\section*{Exercise 4.4.7}

\noindent\textbf{Question.}
Let $(X_n)$ be a martingale with $X_0=0$ and $\E X_n^2<\infty$. Show
that
\[
  \Prob\left(\max_{1\le m\le n}X_m\ge\lambda\right)
  \le
  \frac{\E X_n^2}{\E X_n^2+\lambda^2}.
\]

\Knowledge{
\path{durrett_ch4_doob_maximal_inequality};
\path{durrett_ch4_convex_transform_submartingale};
\path{durrett_ch4_orthogonal_martingale_increments}.
}

\noindent\textbf{Solution.}
For any $c>0$, the process
\[
  Y_m=(X_m+c)^2
\]
is a nonnegative submartingale, since it is a convex function of a
martingale. If
\[
  A=\left\{\max_{1\le m\le n}X_m\ge\lambda\right\},
\]
then on $A$ the maximum of $Y_m$ is at least $(\lambda+c)^2$. Doob's
inequality gives
\[
  \Prob(A)
  \le
  \frac{\E Y_n}{(\lambda+c)^2}
  =
  \frac{\E X_n^2+c^2}{(\lambda+c)^2},
\]
because $\E X_n=0$.
Let $v=\E X_n^2$. Optimizing over $c>0$ gives $c=v/\lambda$, and hence
\[
  \Prob(A)\le \frac{v}{v+\lambda^2}.
\]

\section*{Exercise 4.4.8}

\noindent\textbf{Question.}
Let $(X_n)$ be a submartingale and
\[
  \bar X_n=\max_{0\le m\le n}X_m^+,\qquad
  \log^+x=\max(\log x,0).
\]
Prove
\[
  \E\bar X_n
  \le
  (1-e^{-1})^{-1}
  \left\{1+\E\left[X_n^+\log^+(X_n^+)\right]\right\}.
\]

\Knowledge{
\path{durrett_ch4_doob_maximal_inequality};
\path{durrett_ch4_conditional_expectation_properties}.
}

\noindent\textbf{Solution.}
Let $M>0$. Using the tail integral formula and the trivial bound
$\Prob(\bar X_n\wedge M\ge\lambda)\le1$ for $\lambda\le1$,
\[
\begin{aligned}
  \E(\bar X_n\wedge M)
  &=
  \int_0^\infty \Prob(\bar X_n\wedge M\ge\lambda)\,d\lambda\\
  &\le
  1+\int_1^\infty \Prob(\bar X_n\wedge M\ge\lambda)\,d\lambda.
\end{aligned}
\]
For $\lambda>1$, Doob's inequality gives
\[
  \Prob(\bar X_n\wedge M\ge\lambda)
  \le
  \lambda^{-1}\E\left[X_n^+;\bar X_n\wedge M\ge\lambda\right].
\]
Therefore, by Fubini,
\[
\begin{aligned}
  \E(\bar X_n\wedge M)
  &\le
  1+
  \int X_n^+
  \int_1^{\bar X_n\wedge M}\frac{d\lambda}{\lambda}\,d\Prob\\
  &\le
  1+\E\left[X_n^+\log(\bar X_n\wedge M)\right].
\end{aligned}
\]
The calculus inequality
\[
  a\log b\le a\log a+b/e\le a\log^+a+b/e,
  \qquad a,b\ge0,\ b>0,
\]
follows by maximizing $a\log b-b/e$ as a function of $b$. Applying it
with $a=X_n^+$ and $b=\bar X_n\wedge M$ yields
\[
  \E(\bar X_n\wedge M)
  \le
  1+\E[X_n^+\log^+(X_n^+)]+e^{-1}\E(\bar X_n\wedge M).
\]
Rearranging,
\[
  \E(\bar X_n\wedge M)
  \le
  (1-e^{-1})^{-1}
  \left\{1+\E[X_n^+\log^+(X_n^+)]\right\}.
\]
Letting $M\to\infty$ and using monotone convergence proves the result.

\section*{Exercise 4.4.9}

\noindent\textbf{Question.}
Let $(X_n)$ and $(Y_n)$ be square-integrable martingales. Prove
\[
  \E X_nY_n-\E X_0Y_0
  =
  \sum_{m=1}^n
  \E\bigl[(X_m-X_{m-1})(Y_m-Y_{m-1})\bigr].
\]

\Knowledge{
\path{durrett_ch4_orthogonal_martingale_increments};
\path{durrett_ch4_conditional_variance_formula_martingales}.
}

\noindent\textbf{Solution.}
Write
\[
  \Delta X_m=X_m-X_{m-1},\qquad
  \Delta Y_m=Y_m-Y_{m-1}.
\]
Then
\[
  X_n=X_0+\sum_{m=1}^n\Delta X_m,\qquad
  Y_n=Y_0+\sum_{m=1}^n\Delta Y_m.
\]
When $i<j$, $\Delta X_i$ is $\F_{j-1}$-measurable and
\[
  \E(\Delta X_i\Delta Y_j)
  =
  \E\{\Delta X_i\E(\Delta Y_j\mid\F_{j-1})\}=0.
\]
Similarly, terms with $j<i$ vanish, and the terms involving
$X_0$ or $Y_0$ with later martingale increments vanish. Thus only the
matching increments remain:
\[
  \E X_nY_n
  =
  \E X_0Y_0+
  \sum_{m=1}^n\E(\Delta X_m\Delta Y_m).
\]

\section*{Exercise 4.4.10}

\noindent\textbf{Question.}
Let $(X_n)$ be a martingale and put
\[
  \xi_n=X_n-X_{n-1},\qquad n\ge1.
\]
If $\E X_0^2<\infty$ and
\[
  \sum_{m=1}^\infty \E\xi_m^2<\infty,
\]
show that $X_n$ converges almost surely and in $L^2$.

\Knowledge{
\path{durrett_ch4_orthogonal_martingale_increments};
\path{durrett_ch4_lp_martingale_convergence};
\path{durrett_ch4_conditional_variance_formula_martingales}.
}

\noindent\textbf{Solution.}
By orthogonality of martingale increments,
\[
  \E X_n^2
  =
  \E X_0^2+\sum_{m=1}^n\E\xi_m^2.
\]
The right side is bounded uniformly in $n$ by assumption. Therefore
\[
  \sup_n \E|X_n|^2<\infty.
\]
The $L^p$ martingale convergence theorem with $p=2$ implies that
$X_n$ converges almost surely and in $L^2$ to some square-integrable
limit $X_\infty$.

\section*{Exercise 4.4.11}

\noindent\textbf{Question.}
Continue with the notation of Exercise 4.4.10. If $b_m\uparrow\infty$
and
\[
  \sum_{m=1}^\infty \frac{\E\xi_m^2}{b_m^2}<\infty,
\]
show that
\[
  \frac{X_n}{b_n}\to0\qquad\text{a.s.}
\]
In particular, if $\E\xi_n^2\le K<\infty$ and
$\sum_m b_m^{-2}<\infty$, then $X_n/b_n\to0$ a.s.

\Knowledge{
\path{durrett_ch4_orthogonal_martingale_increments};
\path{durrett_ch4_lp_martingale_convergence};
\path{durrett_ch4_martingale_strong_law_from_variances}.
}

\noindent\textbf{Solution.}
Define
\[
  M_n=\sum_{m=1}^n \frac{\xi_m}{b_m}.
\]
This is a martingale, and its increments have square expectations
$\E\xi_m^2/b_m^2$. By Exercise 4.4.10, $(M_n)$ converges almost surely.
Kronecker's lemma now applies to the convergent series
\[
  \sum_{m=1}^\infty \frac{\xi_m}{b_m}.
\]
It gives
\[
  \frac{1}{b_n}\sum_{m=1}^n\xi_m\to0
  \qquad\text{a.s.}
\]
Since
\[
  \frac{X_n}{b_n}
  =
  \frac{X_0}{b_n}
  +
  \frac{1}{b_n}\sum_{m=1}^n\xi_m
\]
and $b_n\to\infty$, we conclude that $X_n/b_n\to0$ almost surely.
The final assertion follows immediately from
\[
  \sum_m \frac{\E\xi_m^2}{b_m^2}
  \le
  K\sum_m b_m^{-2}<\infty.
\]

\section*{Exercise 4.6.1}

\noindent\textbf{Question.}
Let $Z_1,Z_2,\ldots$ be i.i.d. with $\E|Z_i|<\infty$, let $\theta$ be
an independent random variable with finite mean, and set
\[
  Y_i=Z_i+\theta.
\]
Show that
\[
  \E(\theta\mid Y_1,\ldots,Y_n)\to \theta\qquad\text{a.s.}
\]

\Knowledge{
\path{durrett_ch4_levy_upward_theorem};
\path{durrett_ch4_closed_martingales_ui_equivalence};
\path{durrett_ch4_slln_via_backward_martingales}.
}

\noindent\textbf{Solution.}
Let
\[
  \F_n=\sigma(Y_1,\ldots,Y_n),\qquad
  \F_\infty=\sigma(Y_1,Y_2,\ldots).
\]
By the strong law of large numbers,
\[
  \frac1n\sum_{i=1}^nY_i
  =
  \theta+\frac1n\sum_{i=1}^nZ_i
  \to
  \theta+\E Z_1
  \qquad\text{a.s.}
\]
The left side is $\F_\infty$-measurable, so $\theta$ is
$\F_\infty$-measurable. Hence
\[
  \E(\theta\mid\F_\infty)=\theta.
\]
Since $\F_n\uparrow\F_\infty$, Levy's upward theorem gives
\[
  \E(\theta\mid Y_1,\ldots,Y_n)
  =
  \E(\theta\mid\F_n)
  \to
  \E(\theta\mid\F_\infty)
  =
  \theta
  \qquad\text{a.s.}
\]

\section*{Exercise 4.6.2}

\noindent\textbf{Question.}
Let $\Omega=[0,1)$,
\[
  I_{k,n}=[k2^{-n},(k+1)2^{-n}),\qquad 0\le k<2^n,
\]
and $\F_n=\sigma(I_{k,n}:0\le k<2^n)$. Suppose $f$ is Lipschitz on
$[0,1)$, so $|f(t)-f(s)|\le K|t-s|$. Define
\[
  X_n(\omega)
  =
  \frac{f((k+1)2^{-n})-f(k2^{-n})}{2^{-n}}
  \qquad\text{for }\omega\in I_{k,n}.
\]
Show that $(X_n)$ is a martingale, that $X_n\to X_\infty$ a.s. and in
$L^1$, and that
\[
  f(b)-f(a)=\int_a^b X_\infty(\omega)\,d\omega.
\]

\Knowledge{
\path{durrett_ch4_martingale_definition};
\path{durrett_ch4_closed_martingales_ui_equivalence};
\path{durrett_ch4_levy_upward_theorem}.
}

\noindent\textbf{Solution.}
The random variable $X_n$ is $\F_n$-measurable and $|X_n|\le K$. On a
dyadic parent interval
\[
  I=[k2^{-n},(k+1)2^{-n}),
\]
the two children at level $n+1$ have slopes whose average is
\[
\begin{aligned}
  \frac12&\frac{f((2k+1)2^{-(n+1)})-f(2k2^{-(n+1)})}{2^{-(n+1)}}\\
  &+
  \frac12\frac{f((2k+2)2^{-(n+1)})-f((2k+1)2^{-(n+1)})}{2^{-(n+1)}}\\
  &=
  \frac{f((k+1)2^{-n})-f(k2^{-n})}{2^{-n}}.
\end{aligned}
\]
Thus $\E(X_{n+1}\mid\F_n)=X_n$, so $(X_n)$ is a martingale.

Since $|X_n|\le K$, the martingale is uniformly integrable. The
martingale convergence theorem for uniformly integrable martingales
implies
\[
  X_n\to X_\infty\qquad\text{a.s. and in }L^1.
\]
For dyadic $a<b$, the definition of $X_n$ gives, for all sufficiently
large $n$,
\[
  \int_a^b X_n(\omega)\,d\omega=f(b)-f(a).
\]
Letting $n\to\infty$ and using $L^1$ convergence gives the formula for
dyadic $a,b$. For general $a<b$, choose dyadic $a_m\downarrow a$ and
$b_m\uparrow b$ with $a_m,b_m\in[0,1)$. The Lipschitz continuity of
$f$ and integrability of $X_\infty$ allow passage to the limit, giving
\[
  f(b)-f(a)=\int_a^bX_\infty(\omega)\,d\omega.
\]

\section*{Exercise 4.6.3}

\noindent\textbf{Question.}
With the dyadic sigma-fields from Exercise 4.6.2, suppose $f$ is
integrable on $[0,1)$. Show that $\E(f\mid\F_n)$ is a step function and
\[
  \E(f\mid\F_n)\to f\qquad\text{in }L^1.
\]
Conclude that every integrable $f$ can be approximated in $L^1$ by a
step function.

\Knowledge{
\path{durrett_ch4_levy_upward_theorem};
\path{durrett_ch4_partition_conditioning_formula};
\path{durrett_ch4_ui_plus_probability_equals_l1}.
}

\noindent\textbf{Solution.}
On each dyadic atom $I_{k,n}$,
\[
  \E(f\mid\F_n)
  =
  2^n\int_{I_{k,n}}f(x)\,dx.
\]
Thus $\E(f\mid\F_n)$ is a dyadic step function. The dyadic sigma-fields
increase to the Borel sigma-field on $[0,1)$, completed if necessary.
By Levy's upward theorem,
\[
  \E(f\mid\F_n)\to \E(f\mid\F_\infty)=f
  \qquad\text{a.s. and in }L^1.
\]
Therefore, given $\varepsilon>0$, choose $n$ large enough so that
\[
  \int_0^1|\E(f\mid\F_n)-f|\,dx<\varepsilon.
\]
The function $\E(f\mid\F_n)$ is the desired step-function
approximation.

\section*{Exercise 4.6.4}

\noindent\textbf{Question.}
Let $X_n$ take values in $[0,\infty)$, and set
\[
  D=\{X_n=0\text{ for some }n\ge1\}.
\]
Assume that for every $x<\infty$ there is a number $\delta(x)>0$ such
that
\[
  \Prob(D\mid X_1,\ldots,X_n)\ge \delta(x)
  \qquad\text{a.s. on }\{X_n\le x\}.
\]
Use Levy's 0-1 law to prove
\[
  \Prob\left(D\cup\{\lim_{n\to\infty}X_n=\infty\}\right)=1.
\]

\Knowledge{
\path{durrett_ch4_levy_zero_one_law};
\path{durrett_ch4_levy_upward_theorem};
\path{durrett_ch4_conditional_expectation_definition}.
}

\noindent\textbf{Solution.}
Let $\F_n=\sigma(X_1,\ldots,X_n)$ and
$\F_\infty=\sigma(X_1,X_2,\ldots)$. Since $D\in\F_\infty$, Levy's
0-1 law gives
\[
  \Prob(D\mid\F_n)=\E(\1_D\mid\F_n)\to \1_D
  \qquad\text{a.s.}
\]
Now suppose that $D$ does not occur and that $X_n$ does not tend to
$\infty$. Then for some finite $x$ there are infinitely many $n$ with
$X_n\le x$. Along this subsequence,
\[
  \Prob(D\mid\F_n)\ge\delta(x)>0.
\]
Taking the subsequential limit gives $\1_D\ge\delta(x)$, impossible on
$D^c$. Hence, outside a null set, $D^c$ implies $X_n\to\infty$. This is
equivalent to
\[
  \Prob\left(D\cup\{\lim_nX_n=\infty\}\right)=1.
\]

\section*{Exercise 4.6.5}

\noindent\textbf{Question.}
Let $(Z_n)$ be a branching process with offspring distribution
$(p_k)$, and assume $p_0>0$. Use Exercise 4.6.4 to show that
\[
  \Prob\left(\lim_{n\to\infty}Z_n=0\text{ or }\infty\right)=1.
\]

\Knowledge{
\path{durrett_ch4_branching_process_normalized_martingale};
\path{durrett_ch4_levy_zero_one_law}.
}

\noindent\textbf{Solution.}
Let
\[
  D=\{Z_n=0\text{ for some }n\}.
\]
Once a branching process hits $0$, it remains at $0$. If $Z_n=k\le x$,
then all $k$ individuals in generation $n$ can have no children in the
next generation, an event of conditional probability at least $p_0^k$.
Thus, on $\{Z_n\le x\}$,
\[
  \Prob(D\mid Z_0,\ldots,Z_n)\ge p_0^{\lfloor x\rfloor}>0.
\]
Exercise 4.6.4 gives
\[
  \Prob\left(D\cup\{\lim_n Z_n=\infty\}\right)=1.
\]
On $D$, $Z_n=0$ for all sufficiently large $n$, so $\lim_n Z_n=0$.
Therefore
\[
  \Prob\left(\lim_n Z_n=0\text{ or }\infty\right)=1.
\]

\section*{Exercise 4.6.6}

\noindent\textbf{Question.}
Let $X_n\in[0,1]$ be adapted to $\F_n$. Let $\alpha,\beta>0$ with
$\alpha+\beta=1$, and suppose
\[
  \Prob(X_{n+1}=\alpha+\beta X_n\mid\F_n)=X_n,
  \qquad
  \Prob(X_{n+1}=\beta X_n\mid\F_n)=1-X_n.
\]
Show that
\[
  \Prob(\lim_nX_n=0\text{ or }1)=1,
\]
and if $X_0=\theta$, then
\[
  \Prob(\lim_nX_n=1)=\theta.
\]

\Knowledge{
\path{durrett_ch4_martingale_definition};
\path{durrett_ch4_closed_martingales_ui_equivalence};
\path{durrett_ch4_conditional_variance_formula_martingales}.
}

\noindent\textbf{Solution.}
First compute the conditional mean:
\[
\begin{aligned}
  \E(X_{n+1}\mid\F_n)
  &=
  X_n(\alpha+\beta X_n)+(1-X_n)\beta X_n\\
  &=
  \alpha X_n+\beta X_n
  =
  X_n.
\end{aligned}
\]
Thus $(X_n)$ is a bounded martingale, so $X_n\to X_\infty$ a.s. and in
$L^2$.

The increments are
\[
  X_{n+1}-X_n=
  \begin{cases}
  \alpha(1-X_n), & \text{with conditional probability }X_n,\\
  -\alpha X_n, & \text{with conditional probability }1-X_n.
  \end{cases}
\]
Hence
\[
  \E\bigl((X_{n+1}-X_n)^2\mid\F_n\bigr)
  =
  \alpha^2 X_n(1-X_n).
\]
Since $X_n$ converges in $L^2$, the left side has expectations tending
to $0$, and bounded convergence gives
\[
  0=\alpha^2\E[X_\infty(1-X_\infty)].
\]
Thus $X_\infty(1-X_\infty)=0$ a.s., so
$X_\infty\in\{0,1\}$ a.s.

Finally, $\E X_n=\E X_0=\theta$ for all $n$, and bounded convergence
gives
\[
  \theta=\E X_\infty=\Prob(X_\infty=1).
\]

\section*{Exercise 4.6.7}

\noindent\textbf{Question.}
Show that if $\F_n\uparrow\F_\infty$ and $Y_n\to Y$ in $L^1$, then
\[
  \E(Y_n\mid\F_n)\to \E(Y\mid\F_\infty)
  \qquad\text{in }L^1.
\]

\Knowledge{
\path{durrett_ch4_levy_upward_theorem};
\path{durrett_ch4_conditional_jensen_contraction};
\path{durrett_ch4_ui_plus_probability_equals_l1}.
}

\noindent\textbf{Solution.}
By the $L^1$ contraction property of conditional expectation,
\[
\begin{aligned}
  \E\left|\E(Y_n\mid\F_n)-\E(Y\mid\F_n)\right|
  &\le
  \E\left(\E(|Y_n-Y|\mid\F_n)\right)\\
  &=
  \E|Y_n-Y|
  \to0.
\end{aligned}
\]
By Levy's upward theorem,
\[
  \E(Y\mid\F_n)\to\E(Y\mid\F_\infty)
  \qquad\text{in }L^1.
\]
The triangle inequality combines the two limits and proves
\[
  \E(Y_n\mid\F_n)\to\E(Y\mid\F_\infty)
  \qquad\text{in }L^1.
\]

\section*{Exercise 4.7.1}

\noindent\textbf{Question.}
Show that if a backwards martingale has $X_0\in L^p$, then the
convergence in Theorem 4.7.1 occurs in $L^p$.

\Knowledge{
\path{durrett_ch4_backward_martingale_convergence};
\path{durrett_ch4_levy_downward_theorem};
\path{durrett_ch4_conditional_jensen_contraction};
\path{durrett_ch4_uniform_integrability_definition}.
}

\noindent\textbf{Solution.}
Let $(X_n)_{n\le0}$ be a backwards martingale. Then
\[
  X_n=\E(X_0\mid\F_n),\qquad n\le0.
\]
By Theorem 4.7.1 and Theorem 4.7.2,
\[
  X_n\to X_{-\infty}:=\E(X_0\mid\F_{-\infty})
  \qquad\text{a.s. and in }L^1.
\]
Assume $p\ge1$ and $X_0\in L^p$. Conditional Jensen gives
\[
  |X_n|^p\le \E(|X_0|^p\mid\F_n).
\]
The family on the right is uniformly integrable by the conditional
expectation uniform-integrability theorem, applied to $|X_0|^p\in L^1$.
Also $X_{-\infty}\in L^p$ by the $L^p$ contraction property. Hence
\[
  |X_n-X_{-\infty}|^p
  \le
  2^{p-1}(|X_n|^p+|X_{-\infty}|^p)
\]
is a uniformly integrable family. Since $X_n\to X_{-\infty}$ in
probability, uniform integrability implies
\[
  \E|X_n-X_{-\infty}|^p\to0.
\]
Thus the convergence occurs in $L^p$.

\section*{Exercise 4.7.2}

\noindent\textbf{Question.}
Prove the backwards analogue of Theorem 4.6.10. Suppose
$Y_n\to Y_{-\infty}$ a.s. as $n\to-\infty$ and $|Y_n|\le Z$ for all
$n$, where $\E Z<\infty$. If $\F_n\downarrow\F_{-\infty}$ as
$n\to-\infty$, prove
\[
  \E(Y_n\mid\F_n)\to \E(Y_{-\infty}\mid\F_{-\infty})
  \qquad\text{a.s.}
\]

\Knowledge{
\path{durrett_ch4_levy_downward_theorem};
\path{durrett_ch4_conditional_expectation_properties};
\path{durrett_ch4_conditional_jensen_contraction}.
}

\noindent\textbf{Solution.}
For $N\ge1$, define
\[
  W_N=\sup_{m\le -N}|Y_m-Y_{-\infty}|.
\]
Then $0\le W_N\le 2Z$ and $W_N\downarrow0$ a.s. as $N\to\infty$.
For $n\le -N$,
\[
\begin{aligned}
  \left|\E(Y_n\mid\F_n)-\E(Y_{-\infty}\mid\F_n)\right|
  &\le
  \E(|Y_n-Y_{-\infty}|\mid\F_n)\\
  &\le
  \E(W_N\mid\F_n).
\end{aligned}
\]
By Levy's downward theorem,
\[
  \E(W_N\mid\F_n)\to \E(W_N\mid\F_{-\infty})
  \qquad\text{a.s. as }n\to-\infty.
\]
Since $W_N\downarrow0$ and $W_N\le2Z$, conditional monotone convergence
gives
\[
  \E(W_N\mid\F_{-\infty})\downarrow0.
\]
It follows that
\[
  \E(Y_n\mid\F_n)-\E(Y_{-\infty}\mid\F_n)\to0
  \qquad\text{a.s.}
\]
A second application of Levy's downward theorem gives
\[
  \E(Y_{-\infty}\mid\F_n)
  \to
  \E(Y_{-\infty}\mid\F_{-\infty})
  \qquad\text{a.s.}
\]
The triangle inequality proves the claim.

\section*{Exercise 4.7.3}

\noindent\textbf{Question.}
Prove directly from the definition that if
$X_1,X_2,\ldots\in\{0,1\}$ are exchangeable and
$S_n=X_1+\cdots+X_n$, then
\[
  \Prob(X_1=1,\ldots,X_k=1\mid S_n=m)
  =
  \frac{\binom{n-k}{\,n-m\,}}{\binom{n}{m}}.
\]

\Knowledge{
\path{durrett_ch4_definetti_theorem};
\path{durrett_ch4_backward_martingale_convergence}.
}

\noindent\textbf{Solution.}
Condition on the event $\{S_n=m\}$. By exchangeability, every binary
string of length $n$ with exactly $m$ ones has the same conditional
probability. There are $\binom{n}{m}$ such strings.

The event $\{X_1=\cdots=X_k=1\}$ can occur only if the remaining
$n-k$ coordinates contain $m-k$ additional ones. The number of such
strings is
\[
  \binom{n-k}{m-k}
  =
  \binom{n-k}{n-m}.
\]
Therefore
\[
  \Prob(X_1=1,\ldots,X_k=1\mid S_n=m)
  =
  \frac{\binom{n-k}{m-k}}{\binom{n}{m}}
  =
  \frac{\binom{n-k}{\,n-m\,}}{\binom{n}{m}}.
\]

\section*{Exercise 4.7.4}

\noindent\textbf{Question.}
If $X_1,X_2,\ldots\in\R$ are exchangeable and $\E X_i^2<\infty$, prove
\[
  \E(X_1X_2)\ge0.
\]

\Knowledge{
\path{durrett_ch4_definetti_theorem};
\path{durrett_ch4_orthogonal_martingale_increments}.
}

\noindent\textbf{Solution.}
For each $n$,
\[
  0\le \E\left(\sum_{i=1}^n X_i\right)^2.
\]
By exchangeability,
\[
  \E\left(\sum_{i=1}^n X_i\right)^2
  =
  n\E X_1^2+n(n-1)\E(X_1X_2).
\]
Thus
\[
  \E(X_1X_2)\ge -\frac{\E X_1^2}{n-1}
  \qquad\text{for every }n\ge2.
\]
Letting $n\to\infty$ gives $\E(X_1X_2)\ge0$.

\section*{Exercise 4.7.5}

\noindent\textbf{Question.}
Use the first few lines of the proof of Lemma 4.7.7 to conclude that if
$X_1,X_2,\ldots$ are i.i.d. with $\E X_i=\mu$ and
$\Var(X_i)=\sigma^2<\infty$, then
\[
  \binom{n}{2}^{-1}
  \sum_{1\le i<j\le n}(X_i-X_j)^2
  \to 2\sigma^2
  \qquad\text{a.s.}
\]

\Knowledge{
\path{durrett_ch4_backward_martingale_convergence};
\path{durrett_ch4_slln_via_backward_martingales};
\path{durrett_ch4_hewitt_savage_zero_one}.
}

\noindent\textbf{Solution.}
Let
\[
  U_n=\binom{n}{2}^{-1}
  \sum_{1\le i<j\le n}(X_i-X_j)^2.
\]
This is the unordered version of the symmetric average in Lemma 4.7.7
with $\varphi(x,y)=(x-y)^2$. Since $\E\varphi(X_1,X_2)<\infty$, the
bounded case in Lemma 4.7.7 applies first to
$\varphi_M=\varphi\wedge M$, and then monotone/truncation estimates let
$M\to\infty$. Hence
\[
  U_n\to \E[(X_1-X_2)^2]\qquad\text{a.s.}
\]
For independent copies with common mean $\mu$ and variance $\sigma^2$,
\[
  \E[(X_1-X_2)^2]
  =
  \E X_1^2+\E X_2^2-2\E X_1\E X_2
  =
  2(\sigma^2+\mu^2)-2\mu^2
  =
  2\sigma^2.
\]
Therefore $U_n\to2\sigma^2$ almost surely.

\section*{Exercise 4.8.1}

\noindent\textbf{Question.}
Generalize Theorem 4.8.2. Suppose $L\le M$ are stopping times and
$(Y_{M\wedge n})$ is a uniformly integrable submartingale. Show that
\[
  \E Y_L\le \E Y_M
  \qquad\text{and}\qquad
  Y_L\le \E(Y_M\mid\F_L).
\]

\Knowledge{
\path{durrett_ch4_optional_stopping_ui};
\path{durrett_ch4_optional_stopping_integrable_terminal};
\path{durrett_ch4_bounded_stopping_submartingale_expectation}.
}

\noindent\textbf{Solution.}
For bounded truncations, $L\wedge n\le M\wedge n$, so bounded optional
sampling gives
\[
  \E Y_{L\wedge n}\le \E Y_{M\wedge n}.
\]
Uniform integrability of $(Y_{M\wedge n})$ implies
$Y_{M\wedge n}\to Y_M$ in $L^1$; the same stopped-submartingale argument
gives $Y_{L\wedge n}\to Y_L$ in $L^1$. Letting $n\to\infty$ yields
\[
  \E Y_L\le \E Y_M.
\]
For the conditional form, let $A\in\F_L$ and define $R=L$ on $A$ and
$R=M$ on $A^c$. Then $R$ is a stopping time with $L\le R\le M$, so
\[
  \E Y_R\le \E Y_M.
\]
Since $Y_R=Y_L\1_A+Y_M\1_{A^c}$,
\[
  \E(Y_L;A)\le \E(Y_M;A),\qquad A\in\F_L.
\]
This is exactly
\[
  Y_L\le \E(Y_M\mid\F_L)\qquad\text{a.s.}
\]

\section*{Exercise 4.8.2}

\noindent\textbf{Question.}
If $(X_n)$ is a nonnegative supermartingale, prove
\[
  \Prob\left(\sup_{n\ge0}X_n>\lambda\right)\le \frac{\E X_0}{\lambda}.
\]

\Knowledge{
\path{durrett_ch4_nonnegative_supermartingale_convergence};
\path{durrett_ch4_doob_maximal_inequality};
\path{durrett_ch4_stopped_process_supermartingale}.
}

\noindent\textbf{Solution.}
Let $T=\inf\{n:X_n>\lambda\}$. For each $n$, the stopped process
$X_{T\wedge n}$ is a nonnegative supermartingale, so
\[
  \E X_0\ge \E X_{T\wedge n}
  \ge \E(X_T;T\le n)
  \ge \lambda\Prob(T\le n).
\]
Letting $n\to\infty$ gives
\[
  \Prob\left(\sup_{n\ge0}X_n>\lambda\right)=\Prob(T<\infty)
  \le \frac{\E X_0}{\lambda}.
\]

\section*{Exercise 4.8.3}

\noindent\textbf{Question.}
Let $S_n=\xi_1+\cdots+\xi_n$, where the $\xi_i$ are independent with
$\E\xi_i=0$ and $\Var(\xi_i)=\sigma^2$. Let
\[
  T=\min\{n:|S_n|>a\}.
\]
Use the martingale $S_n^2-n\sigma^2$ to show
\[
  \E T\ge a^2/\sigma^2.
\]

\Knowledge{
\path{durrett_ch4_random_walk_martingales};
\path{durrett_ch4_optional_stopping_integrable_terminal};
\path{durrett_ch4_conditional_variance_formula_martingales}.
}

\noindent\textbf{Solution.}
For the bounded stopping time $T\wedge n$,
\[
  \E S_{T\wedge n}^2=\sigma^2\E(T\wedge n).
\]
For a nondegenerate centered random walk, $T<\infty$ almost surely.
Fatou's lemma gives
\[
  \E S_T^2\le \liminf_{n\to\infty}\E S_{T\wedge n}^2=\sigma^2\E T.
\]
Since $|S_T|>a$, we have $\E S_T^2\ge a^2$, so
\[
  \E T\ge a^2/\sigma^2.
\]

\section*{Exercise 4.8.4}

\noindent\textbf{Question.}
Wald's second equation. Let $S_n=\xi_1+\cdots+\xi_n$, where the
$\xi_i$ are independent with $\E\xi_i=0$ and
$\Var(\xi_i)=\sigma^2$. If $T$ is a stopping time with $\E T<\infty$,
prove
\[
  \E S_T^2=\sigma^2\E T.
\]

\Knowledge{
\path{durrett_ch4_wald_equation};
\path{durrett_ch4_random_walk_martingales};
\path{durrett_ch4_orthogonal_martingale_increments}.
}

\noindent\textbf{Solution.}
For $T\wedge n$,
\[
  \E S_{T\wedge n}^2=\sigma^2\E(T\wedge n).
\]
Since $\E T<\infty$, $\E(T\wedge n)\to\E T$. Also
\[
  S_T-S_{T\wedge n}=\sum_{m=n+1}^{T}\xi_m
\]
on $\{T>n\}$, and martingale-increment orthogonality gives
\[
  \E(S_T-S_{T\wedge n})^2=\sigma^2\E(T-T\wedge n)\to0.
\]
Thus $S_{T\wedge n}\to S_T$ in $L^2$, and
\[
  \E S_T^2=\lim_{n\to\infty}\E S_{T\wedge n}^2=\sigma^2\E T.
\]

\section*{Exercise 4.8.5}

\noindent\textbf{Question.}
Let $P(\xi_i=1)=p$ and $P(\xi_i=-1)=q=1-p$, where $p<1/2$. Let
$S_n=S_0+\xi_1+\cdots+\xi_n$ and
\[
  V_0=\min\{n\ge0:S_n=0\}.
\]
Given $\E_xV_0=x/(1-2p)$, compute $\Var_x(V_0)$.

\Knowledge{
\path{durrett_ch4_asymmetric_gamblers_ruin};
\path{durrett_ch4_wald_equation};
\path{durrett_ch4_random_walk_martingales}.
}

\noindent\textbf{Solution.}
Let $d=q-p=1-2p>0$ and $\tau=V_0$. The centered increment
$Y_i=\xi_i-(p-q)=\xi_i+d$ has variance
\[
  \Var(Y_i)=\Var(\xi_i)=1-(p-q)^2=1-d^2.
\]
Starting from $S_0=x$,
\[
  (S_n-x-(p-q)n)^2-n(1-d^2)
\]
is a martingale. At $\tau$, $S_\tau=0$, so optional stopping gives
\[
  \E(-x+d\tau)^2=(1-d^2)\E\tau.
\]
Using $\E\tau=x/d$,
\[
  d^2\E\tau^2-2xd\E\tau+x^2=(1-d^2)\frac{x}{d},
\]
hence
\[
  \E\tau^2=\frac{x^2}{d^2}+\frac{x(1-d^2)}{d^3}.
\]
Therefore
\[
  \Var_x(V_0)=\frac{x(1-d^2)}{d^3}
  =
  x\,\frac{1-(p-q)^2}{(q-p)^3}.
\]
The answer must be of the form $cx$ because the time to go from $x$ to
$0$ is the sum of $x$ independent copies of the time to go down one
level.

\section*{Exercise 4.8.6}

\noindent\textbf{Question.}
Continue with the setup of Exercise 4.8.5. Find the generating function
of $V_0$.

\Knowledge{
\path{durrett_ch4_asymmetric_gamblers_ruin};
\path{durrett_ch4_random_walk_martingales};
\path{durrett_ch4_optional_stopping_integrable_terminal}.
}

\noindent\textbf{Solution.}
Let
\[
  \phi(\theta)=\E e^{\theta\xi_1}=pe^\theta+qe^{-\theta}.
\]
The exponential martingale gives, for $\theta\le0$,
\[
  e^{\theta x}=\E_x\left[\phi(\theta)^{-V_0}\right].
\]
For $0<s<1$, choose $\theta\le0$ with $\phi(\theta)=1/s$. Writing
$y=e^\theta$, solve
\[
  py+\frac{q}{y}=\frac1s.
\]
The root in $(0,1]$ is
\[
  y=\frac{1-\sqrt{1-4pqs^2}}{2ps}.
\]
Thus
\[
  \E_x(s^{V_0})
  =
  \left(
  \frac{1-\sqrt{1-4pqs^2}}{2ps}
  \right)^x.
\]
It has the form $f(s)^x$ because the descent from $x$ to $0$ is a sum
of $x$ independent one-level descent times.

\section*{Exercise 4.8.7}

\noindent\textbf{Question.}
Let $(S_n)$ be a symmetric simple random walk starting at $0$, and let
\[
  T=\inf\{n:S_n\notin(-a,a)\},
\]
where $a$ is a positive integer. Find constants $b,c$ such that
\[
  Y_n=S_n^4-6nS_n^2+bn^2+cn
\]
is a martingale, and use this to compute $\E T^2$.

\Knowledge{
\path{durrett_ch4_symmetric_gamblers_ruin};
\path{durrett_ch4_random_walk_martingales};
\path{durrett_ch4_optional_stopping_integrable_terminal}.
}

\noindent\textbf{Solution.}
For a symmetric $\pm1$ increment,
\[
  \E(S_n+\xi_{n+1})^2\mid\F_n=S_n^2+1,
  \qquad
  \E(S_n+\xi_{n+1})^4\mid\F_n=S_n^4+6S_n^2+1.
\]
Substitution into $\E(Y_{n+1}\mid\F_n)=Y_n$ gives
\[
  -6+2b=0,\qquad 1-6+b+c=0.
\]
Hence
\[
  b=3,\qquad c=2.
\]
At exit, $S_T=\pm a$. Also $\E T=a^2$. Optional stopping gives
$\E Y_T=Y_0=0$, so
\[
  0=a^4-6a^2\E T+3\E T^2+2\E T.
\]
Therefore
\[
  \E T^2=\frac{5a^4-2a^2}{3}.
\]

\section*{Exercise 4.8.8}

\noindent\textbf{Question.}
Let $S_n=\xi_1+\cdots+\xi_n$ be a random walk. Suppose
\[
  \phi(\theta_0)=\E e^{\theta_0\xi_1}=1
\]
for some $\theta_0<0$, and that $\xi_i$ is not constant. Let
\[
  \tau=\inf\{n:S_n\notin(a,b)\}.
\]
Use $X_n=e^{\theta_0S_n}$ to show
\[
  \E X_\tau=1,\qquad
  \Prob(S_\tau\le a)\le e^{-\theta_0a}.
\]

\Knowledge{
\path{durrett_ch4_insurance_ruin_exponential_bound};
\path{durrett_ch4_random_walk_martingales};
\path{durrett_ch4_optional_stopping_integrable_terminal}.
}

\noindent\textbf{Solution.}
Since $\phi(\theta_0)=1$, $X_n=e^{\theta_0S_n}$ is a nonnegative
martingale with $X_0=1$. Optional stopping for the two-sided exit time
gives
\[
  \E X_\tau=\E X_0=1.
\]
On $\{S_\tau\le a\}$, $\theta_0<0$ implies
\[
  X_\tau=e^{\theta_0S_\tau}\ge e^{\theta_0a}.
\]
Thus
\[
  1=\E X_\tau\ge e^{\theta_0a}\Prob(S_\tau\le a),
\]
or
\[
  \Prob(S_\tau\le a)\le e^{-\theta_0a}.
\]

\section*{Exercise 4.8.9}

\noindent\textbf{Question.}
Continue with Exercise 4.8.8. Suppose the $\xi_i$ are integer-valued,
\[
  \Prob(\xi_i<-1)=0,\qquad \Prob(\xi_i=-1)>0,\qquad \E\xi_i>0.
\]
Let $T=\inf\{n:S_n=a\}$, where $a<0$. Show
\[
  \Prob(T<\infty)=e^{-\theta_0a}.
\]

\Knowledge{
\path{durrett_ch4_insurance_ruin_exponential_bound};
\path{durrett_ch4_asymmetric_gamblers_ruin};
\path{durrett_ch4_random_walk_martingales}.
}

\noindent\textbf{Solution.}
Because downward jumps are never less than $-1$, a lower crossing of
$a$ hits $a$ exactly. Let $\tau_b=T_a\wedge T_b$, where $b>0$.
Optional stopping gives
\[
  1=\E e^{\theta_0S_{\tau_b}}
  =
  e^{\theta_0a}\Prob(T_a<T_b)
  +
  \E(e^{\theta_0S_{T_b}};T_b<T_a).
\]
On $\{T_b<T_a\}$, $S_{T_b}\ge b$, so
$e^{\theta_0S_{T_b}}\le e^{\theta_0b}\to0$. Letting $b\to\infty$,
\[
  1=e^{\theta_0a}\Prob(T_a<\infty),
\]
which proves the formula.

\section*{Exercise 4.8.10}

\noindent\textbf{Question.}
Consider a favorable game in which the payoffs are $-1$, $1$, or $2$,
each with probability $1/3$. Starting with fortune $i$, compute the
probability of ever going broke.

\Knowledge{
\path{durrett_ch4_insurance_ruin_exponential_bound};
\path{durrett_ch4_asymmetric_gamblers_ruin}.
}

\noindent\textbf{Solution.}
Let $r=e^\theta\in(0,1)$. The adjustment equation is
\[
  \E e^{\theta\xi}=1
  \quad\Longleftrightarrow\quad
  \frac{r^{-1}+r+r^2}{3}=1.
\]
Equivalently,
\[
  r^3+r^2-3r+1=(r-1)(r^2+2r-1)=0.
\]
The root in $(0,1)$ is $r=\sqrt2-1$. Starting with fortune $i$, going
broke means the cumulative payoff random walk ever hits $-i$. Hence
\[
  \Prob_i(\text{ever go broke})=(\sqrt2-1)^i.
\]

\section*{Exercise 4.8.11}

\noindent\textbf{Question.}
Let $S_n$ be an insurance company's assets. In year $n$, premiums
totaling $c>0$ are received and claims $\zeta_n$ are paid, where
$\zeta_n$ is normal with mean $\mu$ and variance $\sigma^2$, with
$\mu<c$. If $S_0>0$ is nonrandom, show
\[
  \Prob(\text{ruin})\le \exp\{-2(c-\mu)S_0/\sigma^2\}.
\]

\Knowledge{
\path{durrett_ch4_insurance_ruin_exponential_bound};
\path{durrett_ch4_random_walk_martingales}.
}

\noindent\textbf{Solution.}
The increment $\xi_n=c-\zeta_n$ is normal with mean $m=c-\mu>0$ and
variance $\sigma^2$. Its moment generating function is
\[
  \E e^{\theta\xi_n}
  =
  \exp\left(\theta m+\frac{\sigma^2\theta^2}{2}\right).
\]
The nonzero solution of $\E e^{\theta\xi_n}=1$ is
\[
  \theta_0=-\frac{2m}{\sigma^2}
  =
  -\frac{2(c-\mu)}{\sigma^2}.
\]
Ruin means the accumulated-gain random walk, started from $0$, falls to
$-S_0$ or below. Exercise 4.8.8 with $a=-S_0$ gives
\[
  \Prob(\text{ruin})
  \le
  \exp(\theta_0S_0)
  =
  \exp\{-2(c-\mu)S_0/\sigma^2\}.
\]

\section*{Exercise 4.9.1}

\noindent\textbf{Question.}
Let $a\in S$,
\[
  T_a=\inf\{n\ge1:X_n=a\},
\]
and define
\[
  f_n=P_a(T_a=n),\qquad u_n=P_a(X_n=a).
\]
\begin{enumerate}[label=(\roman*)]
\item Show that, for $n\ge1$,
\[
  u_n=\sum_{m=1}^n f_m u_{n-m}.
\]
\item Let
\[
  u(s)=\sum_{n\ge0}u_ns^n,\qquad
  f(s)=\sum_{n\ge1}f_ns^n.
\]
Show that
\[
  u(s)=\frac{1}{1-f(s)}.
\]
\end{enumerate}

\Knowledge{
\path{durrett_ch4_arcsine_laws_random_walk};
\path{durrett_ch4_reflection_principle};
\path{durrett_ch4_ballot_theorem}.
}

\noindent\textbf{Solution.}
Since the chain starts at $a$, $u_0=P_a(X_0=a)=1$. For $n\ge1$,
decompose the event $\{X_n=a\}$ according to the first positive return
time to $a$:
\[
  \{X_n=a\}
  =
  \bigcup_{m=1}^n \{T_a=m,\ X_n=a\},
\]
where the union is disjoint. By the Markov property at time $T_a$, on
the event $\{T_a=m\}$ the process restarts from $a$, so
\[
  P_a(X_n=a\mid T_a=m)=P_a(X_{n-m}=a)=u_{n-m}.
\]
Therefore
\[
  u_n
  =
  \sum_{m=1}^n P_a(T_a=m)P_a(X_{n-m}=a)
  =
  \sum_{m=1}^n f_m u_{n-m}.
\]

Now multiply by $s^n$ and sum over $n\ge1$. Treating the series as
formal power series, or for $|s|<1$ where the probability generating
functions converge absolutely,
\[
\begin{aligned}
  u(s)-1
  &=
  \sum_{n\ge1}u_ns^n\\
  &=
  \sum_{n\ge1}\sum_{m=1}^n f_mu_{n-m}s^n\\
  &=
  \left(\sum_{m\ge1}f_ms^m\right)
  \left(\sum_{r\ge0}u_rs^r\right)\\
  &=
  f(s)u(s).
\end{aligned}
\]
Hence
\[
  u(s)(1-f(s))=1,
\]
and so
\[
  u(s)=\frac{1}{1-f(s)}.
\]

\end{document}
